%%%%%%%%%%%%%%%%%%%%%%%%%%%%%%%%%%%%%%%%%%%%%%%%%%%%%%%%%%%%%%%%%%%%%
%% Submission manuscript for Computer Methods in Applied Mechanics
%% and Engineering (CMAME, Elsevier), prepared with the official
%% Elsevier "CAS" template bundle (els-cas-templates). The class
%% cas-sc is the single-column manuscript layout; switch to cas-dc
%% for the double-column production layout.
%%%%%%%%%%%%%%%%%%%%%%%%%%%%%%%%%%%%%%%%%%%%%%%%%%%%%%%%%%%%%%%%%%%%%
\documentclass[a4paper,fleqn]{cas-sc}

%%%%%%%%%%%%%%%%%%%%%%%%%%%%%%%%%%%%%%%%%%%%%%%%%%%%%%%%%%%%%%%%%%%%%
%% Bibliography -- CMAME uses numbered references.
%%%%%%%%%%%%%%%%%%%%%%%%%%%%%%%%%%%%%%%%%%%%%%%%%%%%%%%%%%%%%%%%%%%%%
\usepackage[numbers,sort&compress]{natbib}

%%%%%%%%%%%%%%%%%%%%%%%%%%%%%%%%%%%%%%%%%%%%%%%%%%%%%%%%%%%%%%%%%%%%%
%% Math, algorithms, and project-specific packages
%%%%%%%%%%%%%%%%%%%%%%%%%%%%%%%%%%%%%%%%%%%%%%%%%%%%%%%%%%%%%%%%%%%%%
\usepackage{amsmath}
\usepackage{amssymb}
\usepackage{lyu}
\usepackage{algorithm}
\usepackage{algpseudocode}%[noend]
\usepackage{amsthm}\mathchardef\mhyphen="2D

\newtheorem{theorem}{Theorem}
\newtheorem{lemma}[theorem]{Lemma}
\theoremstyle{definition}
\newtheorem{definition}[theorem]{Definition}
\theoremstyle{remark}
\newtheorem{remark}[theorem]{Remark}
\newtheorem{proposition}[theorem]{Proposition}
\newtheorem{assumption}{Assumption}
\newtheorem{corollary}{Corollary}
\graphicspath{ {../},{../figure} } 
\begin{document}

\shorttitle{Equal-weight branching for non-conservative transport--reaction--diffusion equations}
\shortauthors{L. Lyu and H. Lei}

\title[mode=title]{Adaptive equal-weight branching particle methods for non-conservative transport--reaction--diffusion equations}

\author[1,2]{Liyao Lyu}
\cormark[1]
\ead{lyuliyao@math.ucla.edu}
\credit{Conceptualization, Methodology, Software, Formal analysis, Writing -- original draft}

\author[2,3]{Huan Lei}
\ead{leihuan@msu.edu}
\credit{Conceptualization, Supervision, Funding acquisition, Writing -- review \& editing}

\affiliation[1]{organization={Department of Mathematics, University of California, Los Angeles},
            addressline={520 Portola Plaza},
            city={Los Angeles},
            state={CA},
            country={USA}}

\affiliation[2]{organization={Department of Computational Mathematics, Science \& Engineering, Michigan State University},
            addressline={428 S Shaw Lane},
            city={East Lansing},
            state={MI},
            country={USA}}

\affiliation[3]{organization={Department of Statistics \& Probability, Michigan State University},
            addressline={619 Red Cedar Road},
            city={East Lansing},
            state={MI},
            country={USA}}

\cortext[1]{Corresponding author.}

\begin{abstract}
Non-conservative transport--reaction--diffusion equations describe quantities that are transported,
diffused, and locally created or destroyed, as in chemotaxis, reactive transport, and population dynamics.
Particle methods naturally represent transport and diffusion through stochastic differential equations,
but the reaction term changes mass and also needs to be encoded in the particle representation. The direct
realization assigns each particle an exponential weight; this is unbiased, but under localized growth the
represented mass concentrates onto a few particles and the local effective sample size (ESS) collapses
where the solution grows. We develop an \emph{equal-weight stochastic branching} discretization that
converts reaction-induced weight into integer particle count: mass creation becomes particle creation, so
local resolution increases where mass is produced. The method combines an operator splitting of stochastic
transport with a particle-level reaction step, for which we introduce a minimum-variance integer branching
kernel alongside the Poisson birth--death kernel; the weighted update and both kernels reproduce the
reaction multiplier in conditional mean but differ in conditional variance. We prove a continuous-time
convergence result for the branching measure, with Monte-Carlo ($O(N^{-1/2})$) and reconstruction
contributions. On a controlled localized-growth benchmark, branching is more accurate than both raw and systematically resampled weighted particles at the same initial budget, and comparable to ESS-triggered resampling at the final particle count, without any resampling schedule to tune. We then stress-test the scheme on Keller–Segel concentration and a high-dimensional kinetic benchmark. The branching step itself is mesh-free and dimension-agnostic, even though the density reconstruction is not.
\end{abstract}

\begin{highlights}
\item Equal-weight branching encodes non-conservative reaction as integer particle count.
\item A minimum-variance integer kernel lowers variance versus Poisson birth--death.
\item Branching avoids the local effective-sample-size collapse of weighted particles.
\item It outperforms raw and resampled weighted particles at equal cost, with no tuning.
\item The reaction step is mesh-free and dimension-agnostic, tested to 6D and Keller--Segel.
\end{highlights}

\begin{keywords}
particle methods \sep
equal-weight branching \sep
non-conservative transport--reaction--diffusion \sep
weight degeneracy \sep
chemotaxis
\end{keywords}

\maketitle

\section{Introduction}
Non-conservative transport--reaction--diffusion (TRD) equations describe a quantity that is simultaneously
transported by a velocity field, spread by diffusion, and locally created or destroyed by a reaction term.
On a domain $\Omega$, a scalar instance reads
\begin{equation}\label{equ:intro_trd}
\partial_t u = \nabla\!\cdot\!\big(a(u)\,u\big) + D\,\Delta u + r(u)\,u,
\end{equation}
where the total mass $\int_\Omega u\,\intd\bx$ is not preserved.
Such models arise in chemotaxis~\cite{keller1970initiation,patlak1953random}, reactive and porous-media transport~\cite{vazquez2006porous},
combustion~\cite{fisher1937wave,williams2018combustion}, and population dynamics~\cite{hundsdorfer2013numerical}. In these case, the solution routinely develops \emph{localized growth} or \emph{concentration}: the solution forms sharp, rapidly amplifying peaks where the
reaction creates mass. Grid-based discretizations, such as  finite difference~\cite{hoff1978stability}, finite
element~\cite{codina1998comparison}, finite volume~\cite{filbet2006finite,chertock2008second}, and
discontinuous Galerkin schemes~\cite{ayuso2009discontinuous,xia2009application},  are highly developed for
these problems, but resolving a moving, sharpening peak on a fixed mesh typically requires adaptive mesh
refinement or carefully designed positivity-preserving limiters, and the achievable resolution is tied to the
local mesh size.

Particle discretizations offer a complementary, mesh-free route: the solution is represented as an empirical
measure carried by particles, and the resolution follows the particles rather than a fixed grid. Transport and
diffusion are represented naturally, since a drift--diffusion operator is exactly the generator of a stochastic
differential equation (SDE), and particles can be advanced by integrating it. This is the basis of
particle-in-cell solvers for kinetic equations~\cite{birdsall2018plasma,hockney2021computer}, blob and
regularized-kernel methods~\cite{craig2016blob,carrillo2019blob,huang2017error}, and semi-Lagrangian particle
schemes~\cite{qiu2010conservative,qiu2011conservative,qiu2011positivity}, whose convergence to the continuum
equation is justified through the mean-field/propagation-of-chaos
viewpoint~\cite{fournier2004microscopic,hauray2009wasserstein,chaintron2021propagation}. These developments,
however, are largely confined to \emph{conservative} dynamics, where particle number and total mass are fixed.

However, applying this idea to non-conservative equations raises an essential difficulty: the reaction term changes the total and local mass, and this change must be encoded in the particle representation. The most direct construction keeps the positions fixed under reaction and carries a time-dependent \emph{weight} $w_i$ on each particle, updated as $w_i \leftarrow w_i\,\exp(r_i\tau)$ over a step of size $\tau$. While reweighting itself is standard in sequential Monte Carlo and related particle systems~\cite{doucet2000sequential,liu1998sequential}, as a TRD discretization, it inherits a well-known failure mode: in a region of sustained localized growth, the multipliers $\exp(r_i\tau)$ increase multiplicatively, so the represented mass concentrates onto a small number of large-weight particles. The effective sample size (ESS) then collapses, and it does so precisely in the region where the solution is largest and most needs to be resolved. Increasing the global particle count alleviates but does not fully resolve this local degeneracy; it largely postpones rather than prevents it. The standard remedy is resampling, which duplicates high-weight particles and
resets equal weights; it is effective, and we quantify it as a baseline in our numerical experiment.
But resampling is a global operation that must be scheduled (a trigger threshold or period), it redistributes a
fixed particle budget rather than growing resolution with the created mass, it injects variance of its own when
applied too often, and for a finite (non-probability) measure it requires explicit mass bookkeeping. These are
the gaps that an equal-weight, per-step, local representation of the reaction can close.

Rather than repair the weighted representation, we take a different view of the reaction term. Instead of recording mass creation as a growing scalar weight, we record it as a growing \emph{integer particle count}: each particle is replaced by a random, equal-weight number of offspring whose mean equals the reaction multiplier $\exp(r_i\tau)$, so mass creation becomes particle creation and mass destruction becomes particle removal. Every particle keeps the same weight, and the local resolution automatically increases exactly where the reaction amplifies the solution. This is a stochastic branching (birth--death) process~\cite{harris1972nearest,bansaye2015stochastic}, and it is the natural representation here in two senses. Numerically, it keeps equal weights and adapts resolution to the growth, avoiding the local degeneracy of the weighted update. Physically, it mirrors the reaction itself: $r(u)$ creates and destroys mass of the existing field, and branching does the same to particles. The branching kernel is then a design degree of freedom that controls variance and local resolution, and we introduce two such kernels: a Poisson birth--death kernel and a minimum-variance integer kernel, which reproduce the reaction multiplier in conditional mean but differ in conditional variance.


The scheme uses an operator splitting to separate the transport from the reaction: a local Kolmogorov linearization of the drift--diffusion part of \eqref{equ:intro_trd} realizes the transport as an SDE step, and the reaction is applied per particle, so that it needs no global weight normalization and acts directly on a finite (non-probability) measure. The particle viewpoint has a firm theoretical basis for chemotaxis: Stevens derived the Keller--Segel system as the mean-field limit of moderately interacting particles~\cite{Stevens_SIAM_2000}. Whereas that work establishes the continuous mean-field limit, we provide a fully discrete particle scheme -- finite particle number, time stepping, and reconstruction -- for non-conservative TRD equations, of which Keller--Segel is one instance. Related particle methods for these systems~\cite{wang2025novel,hu2024stochastic,chertock2025hybrid} represent transport and diffusion by particles but handle the reaction through a field update, and neural-network solvers based on Feynman--Kac or backward-SDE formulations~\cite{han2024deep,hutzenthaler2021speed,nguwi2024deep,han2018solving,hure2020deep} target a function-space solution; here the reaction is represented at the particle level and the solution is kept as a finite measure throughout. We exercise the scheme on its hardest case, Keller--Segel chemotaxis, whose two-dimensional dynamics can concentrate toward a Dirac profile in finite time~\cite{herrero1997blow,perthame2007transport}, and on a high-dimensional kinetic benchmark.

\paragraph{Contributions}
This paper develops and validates equal-weight branching as a particle discretization for non-conservative
TRD equations. Specifically:
\begin{itemize}
\item[(i)] an equal-weight stochastic branching particle discretization for finite-measure non-conservative
TRD equations, obtained by operator splitting of an SDE transport step and a particle-level reaction step.
\item[(ii)] a unified treatment of the reaction step that places the weighted update, the Poisson
birth--death kernel, and a \emph{minimum-variance integer branching} kernel side by side, with a proposition
showing all three are unbiased for the reaction multiplier while differing in conditional variance.
\item[(iii)] a continuous-time error analysis that proves convergence of the branching measure to the exact
solution, with Monte-Carlo ($O(N^{-1/2})$) and reconstruction contributions and full proofs in
Appendix~\ref{app:convergence}.
\item[(iv)] a controlled numerical comparison showing that branching keeps local
resolution where the weighted scheme degenerates, is more accurate than every weighted baseline at equal
initial budget, and matches ESS-triggered resampling at the matched final particle count.
\item[(v)] application to Keller--Segel concentration -- the hardest, concentration-dominated regime, where the
solution can approach a Dirac profile in finite time -- as a flagship stress test, together with a
high-dimensional kinetic stress test that illustrates the dimension-agnostic nature of the branching step.
\end{itemize}



\section{Equal-weight branching particle discretization}
\label{sec:method}
This section develops the discretization for a scalar non-conservative equation. The construction has four
ingredients: a finite-measure particle representation, an operator splitting into
transport and reaction, an SDE transport step, and a particle-level reaction step in which we compare the weighted
update with two equal-weight branching kernels. The empirical measure is
reconstructed by a truncated Fourier projection, and the extension to coupled
systems, such as Keller--Segel is given in Section~\ref{sec:extension}.

Let us start with the following scalar non-conservative advection-diffusion–reaction system on the $d$-dimensional torus $\mathbb{T}^d$ for a time $t \in (0, T]$:
\begin{subequations}\label{equ:reaction-diffusion}
\begin{align}
\partial_t u &= \nabla \!\cdot\! \big( a(u) u \big) + D \Delta u + r(u) u, & \text{for } (x, t) \in \mathbb{T}^d \times (0, T], \label{equ:pde} \\
u(x, 0) &= u_0(x), & \text{for } x \in \mathbb{T}^d. \label{equ:ic}
\end{align}
\end{subequations}
Here, $u(x,t)$ is the scalar field, $a(u): \mathbb{R}_{+} \to \mathbb{R}^d$ is the advective velocity, $D > 0$ is the diffusion coefficient, and $r(u): \mathbb{R}_{+} \to \mathbb{R}$ is the reaction rate. The system is initialized with $u_0(x)$.
Let $S(N,t)=\{\bX^N_i\}_{i\in\mathcal I(N,t)}$ denote a set of particles at time $t$ whose initial population is $N$, where the particles' indices are numbered consecutively in the index set $\mathcal I(N,t)$ and $\bX^N_i(t)$ denote the position of the $i\mhyphen$th particle. The associated empirical (finite) measure is
\begin{equation}
    t \to \mu^N(t,\intd \bx) =  \frac{1}{N}\sum_{i\in \mathcal I(N,t)}\delta_{\bX^N_i(t)},
\end{equation}
where $\delta_\bx $ denotes the Dirac measure at $\bx\in \mathbb{T}^d$. The core idea of particle-based methods is to approximate the PDE solution by the limit of the empirical measure
\[
\mu(t,\mathrm{d}\bx) = \lim_{N\to\infty} \mu^N(t,\mathrm{d}\bx).
\]
Since Eq. \eqref{equ:reaction-diffusion} is non-conservative, the total mass of the system is not preserved. Consequently, both $\mu^N$ and its limit $\mu$ should be viewed as finite measures rather than probability measures.


\subsection{Initial condition and normalization}
To connect the solution $u(t,\bx)$ to the measure $\mu(t,\mathrm{d}\bx)=\rho(t,\bx) \intd \bx$, we introduce a linear transformation between them. 
We assume throughout that the initial data $u_0$ is nonnegative with positive finite total mass, so that $\mu(0,\cdot)$ can be normalized to a probability measure of total mass $1$; we define
\begin{equation*}
    u(t,\bx)\,\mathrm{d}\bx = Z \,\mu(t,\mathrm{d}\bx), \quad 
Z = \int_{\Omega} u(0,\bx)\,\mathrm{d}\bx
\end{equation*}
denotes the total initial mass.
Thus $\rho(0,\cdot)$ is a probability density
\[
\rho(0,\mathbf{x}) \;=\; \frac{u(0,\mathbf{x})}{Z},
\]
and we can sample the initial particle positions independently from $\rho(0,\cdot)$ (e.g., via Metropolis–Hastings) to form $S(N,0)=\{\mathbf{X}^i_0\}_{i=1}^N$. 
Substituting $u=Z\rho$ into~\eqref{equ:reaction-diffusion}  yields
\begin{equation}\label{equ:normalized_diffusion_reaction}
\partial_t \rho 
= \nabla \cdot \left( \tilde a(\rho)\rho \right)+ D \Delta\rho 
  + \tilde r(\rho)\rho,
\end{equation}
with rescaled coefficients
\[
\tilde a(\rho) = a(Z\rho),  
\qquad \tilde r(\rho) = r(Z\rho).
\]
Since reactions generally change the total mass, $\int_{\Omega}\rho(t,\mathbf{x})\,\mathrm{d}\mathbf{x}$ may not equal $1$ for $t>0$. Rather, it tracks the mass of $u$ via $\int_{\Omega} u(t,\mathbf{x})\,\mathrm{d}\mathbf{x} = Z \int_{\Omega} \rho(t,\mathbf{x})\,\mathrm{d}\mathbf{x}$.

\subsection{Operator splitting}
\label{sec:splitting}
To approximate the solution of the PDE, we discretize the time interval $[0,T]$ with step size $\tau>0$ and grid points
\[
0 = t_0 < t_1 < \cdots < t_M = T,\qquad \tau=t_{n+1}-t_n.
\]
Then an operator-splitting strategy is used to separate the dynamics according to their physical mechanisms. 
Equation~\eqref{equ:normalized_diffusion_reaction} can be written as
\begin{equation}
\partial_t \rho = \mathcal{A}\rho + \mathcal{B}\rho ,
\end{equation}
where the operators $\mathcal{A}$ and $\mathcal{B}$ correspond to
\begin{align}
\mathcal{A}\rho &= \nabla \cdot \left( \tilde a(\rho)\rho \right) + D\Delta\rho  , \\
\mathcal{B}\rho &= \tilde r(\rho)\rho .
\end{align}
Here, $\mathcal{A}$ describes the advection--diffusion dynamics, while $\mathcal{B}$ accounts for the reaction mechanism (e.g., proliferation or death).  The evolution of the exact solution over $t$ can be expressed formally as
\begin{equation}
\rho(t,\bx) = \exp\left( t (\mathcal{A} + \mathcal{B}) \right)\rho(0,\bx).
\end{equation}
In practice, the operator $\exp\!\big(\tau(\mathcal{A}+\mathcal{B})\big)$ cannot be evaluated explicitly, 
we adopt the first-order Lie--Trotter splitting \cite{trotter1959product}. The Trotter product formula guarantees
convergence of the splitting in the limit, and a Baker--Campbell--Hausdorff expansion gives the first-order local
accuracy in $\tau$ that controls the time-discretization term of the error decomposition~\eqref{eq:err-decomp}:
\begin{equation}
\exp\!\left( t (\mathcal{A} + \mathcal{B}) \right) 
= \lim_{M\to\infty} \Big( \exp\!\left(\tfrac{t}{M} \mathcal{B}\right) 
                          \exp\!\left(\tfrac{t}{M} \mathcal{A}\right) \Big)^M .
\end{equation}
At the discrete level, we approximate $\rho(t_{n+1},\cdot)$ by
\begin{equation}
\rho_{n+1} = \exp(\tau \mathcal{B}) \exp(\tau \mathcal{A}) \rho_n ,
\end{equation}
for $n=0,\cdots,M-1$.
This scheme provides a simple and efficient first-order approximation; 
higher-order splittings, such as Strang splitting \cite{strang1968construction}, may be employed for improved accuracy.
In practice, this decomposition allows us to design particle-based updates for each subprocess separately:
\begin{equation}
\rho_n^\ast = \exp(\tau \mathcal{A})\rho_n, 
\qquad 
\rho_{n+1} = \exp(\tau \mathcal{B})\rho_n^\ast.
\end{equation}

\subsection{Advection--diffusion step}
\label{sec:advection_diffusion}
In this step, we propagate particles under the advection–diffusion dynamics, where each particle follows a stochastic path determined by both a drift and a diffusion term. At time step $n$, let $S^N_n = \{\bX^N_{i,n}\}_{i\in \mathcal I_n^N}$ denote  a set of particles  whose initial population is $N$, which induce the empirical measure 
\[
\mu^N_{n}= \frac{1}{N}\sum_{i\in \mathcal I_n^N}\delta_{\bX^N_{i,n}}
\] 
approximates the solution at step $n$ with a finite number of particles, i.e.,  $\lim_{N\to\infty} \mu^N_{n} (\intd \bx) = \mu_n(\intd \bx) = \rho_n(\bx)\intd \bx$. 
Our goal is to evolve these particles according to the transition probability $p(\bx,s;\by,t)$, which leads to an updated empirical measure
\[\mu_n^{N,*} = \int p(\bx,t_n;\by,t_n+\tau)  \mu_n^N(\intd \bx) =   \frac{1}{N}\sum_{i\in \mathcal I_n^N}\delta_{\bX_{i,n}^{N,*}}
\]  
reflecting the action of the advection--diffusion operator on the density.
By the Kolmogorov forward equation, the transition density $p(\bx,s;\by,t)$ satisfies 
\begin{equation}
    \partial_t p = \tilde {\mathcal A }p := - \nabla \cdot(b  p) + \frac{1}{2}\nabla^2:(\sigma \sigma^\top p) 
\end{equation}
with initial condition $\lim_{t\to s^+} p(\bx,s;\cdot,t) = \delta_\bx$. By matching the differential operator $\mathcal A$ with the target PDE \eqref{equ:reaction-diffusion}, we have 
\begin{equation}
    b(x) = -\tilde a(\rho^N_n(x)), \quad \frac{1}{2}\sigma \sigma^\top = D,
\end{equation}
where $\rho^N_n(x)$ represents the interpolated density of $\mu_n^N$ and will be discussed in Sec. \ref{sec:interpolation}.
Within each step the coefficients are frozen at the current reconstruction: the transport drift $b$ is evaluated
from $\rho^N_n$ and held fixed over $[t_n,t_{n+1}]$, and the reaction rate is then evaluated from the
post-transport reconstruction $\rho^{N,\ast}_n$ in the subsequent reaction substep (Section~\ref{sec:reaction}).
%
Accordingly, the evolution of each particle $\bX_{i}^{N}(t)$ as a stochastic process satisfying the It\^o stochastic differential equation
\begin{equation}\label{equ:overd_langevin}
    \intd \bX^N_{i}(t) = b(\bX^N_{i} (t)) \intd t  + \sigma \intd \bW^N_{i}(t),
\end{equation}
where $\bW^N_i(t)$ is a $d\mhyphen$dimensional Brownian motion.
In practice, we approximate Eq. \eqref{equ:overd_langevin} using an Euler-Maruyama scheme with step size $\tau$, 
\begin{equation}
 \bX_{i,n}^{N,*} = \mathbf{X}_{N,n}^{i} + b(\bX^N_{i,n}) \tau  + \sigma \sqrt{\tau} \boldsymbol{\xi}^i_{N,n},
 \label{eq:euler_maruyama}
\end{equation}
where $\boldsymbol{\xi}^i_{N,n}$ are independent and identically distributed standard normal random variables; see Algorithm \ref{alg:sde} for the summary of the particle propagation for the advection--diffusion step. 

\begin{algorithm}
\caption{SDE propagation of the advection--diffusion step}\label{alg:sde}
\begin{algorithmic}
\Require Particle set: $S=\{\mathbf{X}_i\}_{i\in\mathcal I}$; time step $\tau>0$
\Require Drift term $b(\cdot)$
\Require Diffusion strength $\sigma$
\Ensure Propagated set $S_*=\{\mathbf{X}^{i,*}\}_{i\in\mathcal I}$
\Procedure{SDE}{$S,\,\tau,\,b,\,\sigma$}
  \State $S_{*}\gets [\ ]$
  \For{$i\in\mathcal I$}
    \State Sample $\boldsymbol{\xi}_i \sim \mathcal N(0,I_d)$ independently
    \State $\mathbf{X}^*_i \gets \mathbf{X}_i \;+\; b(\mathbf{X}_i)\,\tau \;+\; \sqrt{\tau}\,\sigma\,\boldsymbol{\xi}_i$
    \State Append $\mathbf{X}^{*}_i$ to $S_{*}$
  \EndFor
  \State \Return $S_{*}$
\EndProcedure
\end{algorithmic}
\end{algorithm}

\subsection{Reaction step: weighted update versus branching}
\label{sec:reaction}
The reaction sub-step advances the intermediate measure $\mu_n^{N,\ast}$ from the transport step under the
reaction operator $\mathcal B$ over a time step $\tau$, that is, it solves the ordinary
differential equation
\[
\partial_t \rho = c_n(\mathbf{x})\,\rho,
\qquad c_n(\mathbf{x}) := \tilde r\!\big(\rho_n^\ast(\mathbf{x})\big),
\]
whose exact solution multiplies the density pointwise by the \emph{reaction multiplier}
\begin{equation}\label{equ:multiplier}
m(\bx) = \exp\!\big(c_n(\bx)\,\tau\big)\ge 0 .
\end{equation}
At the particle level we must therefore produce, from $\mu_n^{N,\ast}=\frac1N\sum_i \delta_{\bX_i^\ast}$, a new
measure $\mu_{n+1}^N$ with
\begin{equation}\label{equ:reaction_target}
\mathbb E\big[\langle f,\mu_{n+1}^N\rangle\big]
=\langle f, m\,\mu_n^{N,\ast}\rangle
=\frac1N\sum_{i} m_i\, f(\bX_i^\ast),
\qquad m_i := m(\bX_i^\ast),
\end{equation}
for all $f\in C_b(\mathbb T^d)$, i.e.\ the update must be \emph{unbiased} for the reaction multiplier. There is
more than one particle realization of \eqref{equ:reaction_target}, and the choice has a decisive effect on the
quality of the discretization. We describe the weighted update and two equal-weight branching kernels, and then
contrast their conditional variance and resolution.

\paragraph{Weighted update (baseline)}
The most direct realization keeps the positions fixed and carries a scalar weight $w_i$ on each particle,
\begin{equation}\label{equ:weighted}
w_i^{n+1} = w_i^{n}\,m_i,
\qquad
\mu_w^{n+1}(\intd\bx) = \frac1N\sum_i w_i^{n+1}\,\delta_{\bX_i^\ast}(\intd\bx),
\end{equation}
which satisfies \eqref{equ:reaction_target} exactly and deterministically (the weight update has zero variance).
Its drawback is that the weights become unequal. A standard measure of the resulting degeneracy is the
normalized effective sample size (nESS),
\begin{equation}\label{equ:ESS}
\mathrm{nESS} = \frac{1}{N_{\mathrm{act}}}\,\frac{\big(\sum_i w_i\big)^2}{\sum_i w_i^2}\in(0,1],
\end{equation}
where $N_{\mathrm{act}}$ is the number of active particles; equal weights give $\mathrm{nESS}=1$, while
$\mathrm{nESS}\to0$ indicates that the represented mass has concentrated onto a few particles. Under sustained
localized growth the multipliers $m_i$ compound multiplicatively in the high-rate region, so $\mathrm{nESS}$
decays there even though the global particle count is unchanged -- the local degeneracy described in the
introduction. By construction $\mathrm{nESS}$ diagnoses the \emph{weighted} representation only: branching keeps
all weights equal, so $\mathrm{nESS}\equiv1$ for the branching kernels. The corresponding resolution diagnostic
for branching is the local particle count $N_B$, which we compare against the weighted local $\mathrm{nESS}(B)$
in Section~\ref{sec:branch_vs_weight}.

\paragraph{Equal-weight branching}
Branching instead keeps all weights equal to $1/N$ and replaces particle $i$ by a random integer number of
identical offspring $\nu_i\in\mathbb Z_{\ge0}$ with $\mathbb E[\nu_i]=m_i$. The updated set
$S_{n+1}^N=\{\bX_{j,n+1}^N\}_{j\in\mathcal I_{n+1}^N}$ collects the offspring, and
\begin{equation}
\mu_{n+1}^{N}(\intd\bx)=\frac1N\sum_{j\in\mathcal I_{n+1}^N}\delta_{\bX_{j,n+1}^N}(\intd\bx),
\end{equation}
so that \eqref{equ:reaction_target} holds because $\mathbb E[\nu_i]=m_i$. Mass creation ($m_i>1$) becomes
particle creation and mass loss ($m_i<1$) becomes removal, so the local particle count -- and hence the local
resolution -- adapts to where the reaction acts, while every particle keeps weight $1/N$. We use two offspring
distributions.

\emph{(i) Poisson birth--death kernel} (a Poisson birth rule for $m_i\ge1$ and a Bernoulli survival rule for
$m_i<1$). For $c_i\ge0$ (so $m_i\ge1$), keep the parent and add
$\mathrm{Poisson}(m_i-1)$ children; for $c_i<0$ (so $m_i\in(0,1)$), keep the parent with probability $m_i$:
\begin{equation}\label{equ:poisson_kernel}
\nu_i =
\begin{cases}
1+\mathrm{Poisson}(m_i-1), & c_i\ge 0,\\[2pt]
\mathrm{Bernoulli}(m_i), & c_i<0,
\end{cases}
\end{equation}
which satisfies $\mathbb E[\nu_i]=m_i$, with $\operatorname{Var}(\nu_i)=m_i-1$ for $c_i\ge0$ and
$\operatorname{Var}(\nu_i)=m_i(1-m_i)$ for $c_i<0$.

\emph{(ii) Minimum-variance integer kernel.} The offspring count must be an integer with mean $m_i$, yet
$m_i=e^{c_i\tau}$ is generally non-integer. We therefore round $m_i$ stochastically to one of its two bracketing
integers $\lfloor m_i\rfloor$ and $\lfloor m_i\rfloor+1$ -- the stochastic-rounding rule, here used as an
offspring distribution:
\begin{equation}\label{equ:minvar_kernel}
\nu_i=\lfloor m_i\rfloor + \mathrm{Bernoulli}(\theta_i),
\qquad
\theta_i = m_i-\lfloor m_i\rfloor,
\end{equation}
which gives $\mathbb E[\nu_i]=m_i$ and $\operatorname{Var}(\nu_i)=\theta_i(1-\theta_i)$. Among all
integer-valued offspring distributions with mean $m_i$, this is the one of smallest variance
(Proposition~\ref{prop:stochastic_branch}); since $\theta_i\in[0,1)$, the variance never exceeds $\tfrac14$ and vanishes
whenever $m_i$ is an integer, where $\nu_i\equiv m_i$ deterministically.

For $m_i\in(0,1)$ the rule reduces to $\mathrm{Bernoulli}(m_i)$, coinciding with the Poisson death rule. For
$m_i>1$ its variance $\theta_i(1-\theta_i)$ is strictly smaller than the Poisson birth variance $m_i-1$, the gap
being largest at integer $m_i$, where the Poisson rule still has variance $m_i-1$ while the present rule is
deterministic.
\begin{proposition}[Unbiasedness and conditional variance of the reaction step]
\label{prop:stochastic_branch}
Let $\mu=\frac1N\sum_{i\in\mathcal I}\delta_{\bX_i}$ and let $\mu_+$ be produced by any of the three reaction
updates of Section~\ref{sec:reaction}, with multipliers $m_i=\exp(c_i\tau)$ and $c_i:=c(\bX_i)$. Then,
conditionally on $\{\bX_i\}_{i\in\mathcal I}$,
\[
\mathbb E[\langle f,\mu_+\rangle]=\langle f,\, m\,\mu\rangle
=\frac1N\sum_{i\in\mathcal I} m_i\,f(\bX_i)
\qquad\text{for all } f\in C_b(\mathbb T^d),
\]
so all three updates are unbiased for the reaction multiplier. The conditional variance of $\langle f,\mu_+\rangle$
is
\[
\operatorname{Var}\!\big(\langle f,\mu_+\rangle \,\big|\, \{\bX_i\}\big)
=\frac1{N^2}\sum_{i\in\mathcal I}\operatorname{Var}(\nu_i)\,f(\bX_i)^2,
\]
with $\operatorname{Var}(\nu_i)$ as in \eqref{equ:poisson_kernel}--\eqref{equ:minvar_kernel} for the two
branching kernels, and is $0$ for the weighted update. Moreover, among all integer-valued offspring distributions
with mean $m_i$, the minimum-variance kernel \eqref{equ:minvar_kernel} attains the smallest possible variance,
$\theta_i(1-\theta_i)$ with $\theta_i=m_i-\lfloor m_i\rfloor$.
\end{proposition}

\begin{proof}
We treat unbiasedness, the conditional variance, and optimality in turn.

\emph{Unbiasedness.} For either branching kernel the $\nu_i$ offspring of particle $i$ are placed at $\bX_i$, so
$\langle f,\mu_+\rangle=\frac1N\sum_{i\in\mathcal I}\nu_i\,f(\bX_i)$. Taking the conditional expectation and using
$\mathbb E[\nu_i\mid\bX_i]=m_i$ gives
$\mathbb E[\langle f,\mu_+\rangle]=\frac1N\sum_i m_i\,f(\bX_i)=\langle f,m\mu\rangle$. For the weighted update
$\langle f,\mu_+\rangle=\frac1N\sum_i w_i m_i\,f(\bX_i)$ holds deterministically, with the same value in
expectation.

\emph{Conditional variance.} The offspring counts $\nu_i$ are independent across $i$ given $\{\bX_i\}$, so
\[
\operatorname{Var}\!\big(\langle f,\mu_+\rangle\big)
=\frac1{N^2}\sum_{i\in\mathcal I}\operatorname{Var}(\nu_i)\,f(\bX_i)^2 .
\]
The single-particle variances are $\operatorname{Var}(\nu_i)=m_i-1$ for the Poisson birth rule
($\nu_i=1+\mathrm{Poisson}(m_i-1)$), $m_i(1-m_i)$ for the Bernoulli death rule, and $\theta_i(1-\theta_i)$ for the
minimum-variance kernel \eqref{equ:minvar_kernel}, where $\nu_i\in\{\lfloor m_i\rfloor,\lfloor m_i\rfloor+1\}$
with probabilities $1-\theta_i,\theta_i$. The weighted update is deterministic, so its variance is $0$.

\emph{Optimality of \eqref{equ:minvar_kernel}.} Fix $i$ and let $\nu$ be any integer-valued random variable with
$\mathbb E[\nu]=m_i$. Write $k=\lfloor m_i\rfloor$ and $\theta=\theta_i=m_i-k\in[0,1)$, and let
$L(n)=\theta^2+(1-2\theta)(n-k)$ be the secant line of the convex map $n\mapsto(n-m_i)^2$ through $n=k$ and
$n=k+1$. Since no integer lies strictly inside $(k,k+1)$, convexity gives $(n-m_i)^2\ge L(n)$ for every integer
$n$, with equality iff $n\in\{k,k+1\}$. Taking expectations and using $\mathbb E[\nu]=m_i=k+\theta$,
\[
\operatorname{Var}(\nu)=\mathbb E[(\nu-m_i)^2]\ge\mathbb E[L(\nu)]
=\theta^2+(1-2\theta)\theta=\theta(1-\theta),
\]
with equality iff $\nu\in\{k,k+1\}$ almost surely. This is precisely the kernel \eqref{equ:minvar_kernel}, which
is therefore the minimum-variance integer-valued unbiased offspring distribution.
\end{proof}
Table~\ref{tab:kernels} summarizes the three reaction representations. All are unbiased; they differ in how the
reaction is stored. The weighted update has zero sampling variance but lets weights -- and therefore the
effective sample size -- degenerate locally. Branching keeps weights equal and trades a controlled offspring
variance for adaptive local resolution, and the minimum-variance kernel minimizes that variance among
integer-valued unbiased kernels. We summarize the equal-weight reaction step in Algorithm~\ref{alg:birthdeath};
the weighted update simply replaces lines~\ref{alg:line:birth}--\ref{alg:line:death} by $w_i\leftarrow w_i m_i$.

\begin{table}[t]
\centering
\caption{Reaction representations of the multiplier $m_i=\exp(c_i\tau)$, with $\theta_i=m_i-\lfloor m_i\rfloor$.
All three are unbiased ($\mathbb E[\nu_i]=m_i$, or an exact factor $m_i$ for weights). ``Local resolution''
refers to whether the number of particles in a growth region increases with the created mass.}
\label{tab:kernels}
\begin{tabular}{lccc}
\hline
Representation & weight & conditional variance & local resolution \\
\hline
Weighted update \eqref{equ:weighted}        & $w_i m_i$ (unequal) & $0$ & fixed count; nESS degenerates \\
Poisson kernel \eqref{equ:poisson_kernel}   & $1/N$ (equal) & $m_i-1$ \,/\, $m_i(1-m_i)$ & adaptive \\
Min.-variance kernel \eqref{equ:minvar_kernel} & $1/N$ (equal) & $\theta_i(1-\theta_i)\le \tfrac14$ & adaptive \\
\hline
\end{tabular}
\end{table}


\begin{algorithm}
\caption{Equal-weight reaction (branching) step}\label{alg:birthdeath}
\begin{algorithmic}[1]
\Require Propagated particles $S=\{\mathbf X_{i}\}_{i\in\mathcal I}$, time step $\tau>0$, rate function $r(\cdot)$, kernel $\in\{\textsc{Poisson},\textsc{MinVar}\}$
\Ensure Post-reaction set $S_{+}=\{\mathbf X_{i}\}_{i\in\mathcal I_{+}}$ (all weights $1/N$)
\Procedure{Reaction}{$S,\,\tau,\,r(\cdot),\,\text{kernel}$}
    \State $S_{+}\gets [\ ]$
    \For{$i\in\mathcal I$}
        \State $m_i \gets \exp\!\big(r(\mathbf X_{i})\,\tau\big)$
        \If{kernel $=\textsc{Poisson}$} \label{alg:line:birth}
            \State $\nu_i \gets 1+\mathrm{Poisson}(m_i-1)$ if $m_i\ge1$, else $\nu_i\gets\mathrm{Bernoulli}(m_i)$
        \Else \Comment{minimum-variance}
            \State $\nu_i \gets \lfloor m_i\rfloor + \mathrm{Bernoulli}(m_i-\lfloor m_i\rfloor)$ \label{alg:line:death}
        \EndIf
        \State Append $\nu_i$ copies of $\mathbf X_{i}$ to $S_{+}$
    \EndFor
    \State \Return $S_{+}$
\EndProcedure
\end{algorithmic}
\end{algorithm}


\subsection{Nonlinear interpolation}
\label{sec:interpolation}
To evaluate the functions $\tilde a(\rho)$ , and $\tilde r(\rho)$ , that govern the drift and reaction dynamics in our particle-based scheme, we need to reconstruct the continuous density 
$\rho$ from its empirical measure. Specifically, given an empirical measure
\[
\mu^N_n= \frac{1}{N}\sum_{i\in \mathcal I^N_n}\delta_{\bX_{i,n}^N},
\] 
we seek a smooth approximation of $\rho^N_n$ that allows pointwise evaluation and differentiation as required by the operator $\mathcal{A}$ and $\mathcal{B}$.
In this work, we choose to interpolate the empirical distribution in the Fourier space. Specifically, we project the density onto a finite set of Fourier basis functions, allowing efficient representation and differentiation. 
The Fourier coefficients are evaluated from the particle locations $\bX^N_{i,n}$ under periodic boundary conditions. Without loss of generality, we take the periodic domain to be $\Omega=[-\pi,\pi]^d$ , from which the reconstructed density $\rho_n(t,x)$. To be more specific, we define a Fourier interpolation operator $\mathcal P_K$
that maps the empirical particle distribution into a truncated Fourier representation 
\[
\rho^N_n(\bx) = \mathcal{P}_K(\mu^N_n)(\bx) 
          = \sum_{\vert k_1\vert ,\dots,\vert k_d\vert \leq K} \hat{\rho}_n(k_1,\dots,k_d)\,\phi_{k_1}(\bx_1)\cdots\phi_{k_d}(\bx_d),
\]
where the one-dimensional basis functions on $[-\pi,\pi]$ are the orthonormal real Fourier modes
\[
\phi_k(x) =
\begin{cases}
\dfrac{1}{\sqrt{2\pi}}, & k=0, \\[4pt]
\dfrac{1}{\sqrt{\pi}} \cos(k  x), & k \in \mathbb{Z}_+, \\[4pt]
\dfrac{1}{\sqrt{\pi}}\sin(|k|  x), & k \in \mathbb{Z}_- ,
\end{cases}
\]
so that $\int_{-\pi}^{\pi}\phi_k\phi_{k'}\,\intd x=\delta_{kk'}$, and the sum runs over all $d$ directions.

The Fourier coefficients are computed directly from the empirical measure, i.e.\ as fixed-$N$ particle averages
of the basis functions:
\begin{equation}
\begin{aligned}
    \hat{\rho}_n(k_1,\dots,k_d) &= \int \phi_{k_1}(y_1)\cdots \phi_{k_d}(y_d)\,\mu^N_n(\intd\by)
    = \frac{1}{N}\sum_{i\in \mathcal I^N_n} \phi_{k_1}(\bX_{i,n,1}^{N})\cdots \phi_{k_d}(\bX_{i,n,d}^{N}).
\end{aligned}
\label{eq:Fourier_mode}
\end{equation}
The process is shown in Algorithm \ref{alg:interpolate}. 
This choice is motivated by the simplicity and spectral accuracy of Fourier methods for smooth periodic functions. However, we note that other functional spaces, such as polynomial bases, wavelets, or kernel-based representations, could offer better adaptability in non-periodic or highly localized settings. We leave the exploration of such alternatives to future work.



\begin{algorithm}
\caption{Interpolation Step (fixed-$N$ normalization)}\label{alg:interpolate}
\begin{algorithmic}
\Require Particle set: $S=\{ \mathbf{X}^i \}_{i\in \mathcal{I}}$; fixed initial population $N$; truncation index set $\mathcal{M}_K$; domain descriptor $\Omega$; basis $\{\psi_m\}_{m\in\mathcal{M}_K}$ with (optionally) known gradients $\{\nabla\psi_m\}_{m\in\mathcal{M}_K}$
\Ensure Truncated field $\tilde{\rho}_n(\mathbf{x})$ (and, if needed, $\nabla\tilde{\rho}_n(\mathbf{x})$)
\Procedure{Interp}{$S_n,\,N,\,\mathcal{M}_K,\,\Omega,\,\{\psi_m\}$}
    \State \textbf{(Coefficient projection with fixed-$N$ normalization)}
    \State For each $m\in\mathcal{M}_K$, set
    \[
      a_m \;\gets\; \frac{1}{N}\sum_{i\in\mathcal{I}} \psi_m\!\big(\mathbf{X}^i\big).
    \]
    \State \textbf{(Series assembly)}
    \[
      \tilde{\rho}_n(\mathbf{x}) \;\gets\; \sum_{m\in\mathcal{M}_K} a_m\,\psi_m(\mathbf{x}).
    \]
    \State \textbf{(Optional gradients for drift/evaluation)}
    \If{gradients are required and $\{\nabla\psi_m\}$ are available}
      \[
        \nabla\tilde{\rho}_n(\mathbf{x}) \;\gets\; \sum_{m\in\mathcal{M}_K} a_m\,\nabla\psi_m(\mathbf{x}).
      \]
    \EndIf
    \State \Return $\tilde{\rho}_n$ (and $\nabla\tilde{\rho}_n$ if computed)
\EndProcedure
\end{algorithmic}
\end{algorithm}

Because the coefficients are normalized by the \emph{initial} population $N_0$ rather than the current count,
the reconstructed field $\tilde\rho_n$ carries the current mass: it integrates to $M_t=N_t/N_0$, i.e.
$\tilde\rho_n = M_t\,\rho_n^{\mathrm{prob}}$ where $\rho_n^{\mathrm{prob}}$ is the unit-mass density of the active
cloud. Consequently $\tilde u = Z\tilde\rho_n$ recovers the physical field with the correct total mass, and the
drift and reaction coefficients are evaluated from this mass-carrying reconstruction throughout (so the factors
$M_t$ enter consistently and do not need separate bookkeeping). In regions where $\tilde\rho_n$ is close to zero
the per-particle rate (for a coupled source $f$, the quotient $\tilde r=f/(Z\tilde\rho_n)=f/\tilde u$) is regularized by a small density floor to keep the multiplier $e^{\tilde r\tau}$ bounded.

\subsection{Optional population control}
\label{sec:popcontrol}
Sustained growth makes the active population $N_t$ grow, which is the mechanism by which branching adds local
resolution but which can also raise the cost. When a cost ceiling is needed, we apply a global, equal-weight
population control that rescales the count without introducing local weights. Writing the current mass as
$M_t=\langle 1,\mu^N\rangle = N_t/N$, we resample $N_{\mathrm{targ}}$ particles uniformly (with replacement)
from the current set and reattach the common mass factor,
\begin{equation}
\mu^N=\frac1N\sum_{i=1}^{N_t}\delta_{\bX_i}
\;\longmapsto\;
M_t\,\frac1{N_{\mathrm{targ}}}\sum_{j=1}^{N_{\mathrm{targ}}}\delta_{\tilde\bX_j},
\end{equation}
which preserves the total mass and the equal-weight property in expectation. This step is optional and is
\emph{not} used in the branching-versus-weighted comparison of Section~\ref{sec:branch_vs_weight}; whenever it is
active we report it explicitly.

\section{Theoretical properties}
\label{sec:theory}
In this section we analyze the error between the numerical solution and the exact
solution of~\eqref{equ:reaction-diffusion}. We decompose the error at a grid time
$t_n$ into three parts:
\begin{equation}\label{eq:err-decomp}
u(t_n)-\mu_n^N
=\underbrace{\bigl(u(t_n)-u_K(t_n)\bigr)}_{\text{coefficient regularization}}
+\underbrace{\bigl(u_K(t_n)-\mu^N(t_n)\bigr)}_{\text{continuous-time particle}}
+\underbrace{\bigl(\mu^N(t_n)-\mu_n^N\bigr)}_{\text{time discretization}}.
\end{equation}
The first term is the consistency error between the exact equation and the
\emph{coefficient-regularized} equation~\eqref{equ:projected}; the latter replaces
only the coefficient arguments $a(u),r(u)$ by $a(\mathcal P_K u_K),r(\mathcal P_K
u_K)$, so it is \emph{not} a spectral Galerkin projection of the equation , and its error is governed by the spectral accuracy
of $\mathcal P_K$ acting on the coefficient argument. The second term is the
continuous-time particle approximation error between $u_K$ and the continuous-time
branching measure $\mu^N$, an additive sampling fluctuation of size $O(1/N)$. The
present paper controls these first two terms and thereby the continuous-time error
$u-\mu^N$ (Theorem~\ref{thm:abstract_error}). The third term is the
time-discretization error between $\mu^N$ and the fully discrete algorithm
$\mu^N_n$; its analysis decomposes the flow by the operator splitting of
Section~\ref{sec:splitting} and compares the continuous- and discrete-time flows
pathwise on a single probability space, through a Harris-type graphical
construction~\cite{harris1972nearest,bansaye2015stochastic} in the
measure-valued/SDE form of~\cite{fournier2004microscopic}. That construction is
carried out in a future work; combined
with Theorem~\ref{thm:abstract_error} it yields the fully discrete bound
\eqref{equ:abstract_rate_full}.

\paragraph{The coefficient-regularized equation.}
We introduce the projected PDE
\begin{equation}\label{equ:projected}
\partial_t u_K = \nabla\cdot\!\left( a(\mathcal P_K u_K)\, u_K\right)
+ D\,\Delta u_K + r\!\left(\mathcal P_K u_K\right) u_K,
\end{equation}
where $\mathcal P_K$ is the Fourier projection of Section~\ref{sec:interpolation}
onto the modes with $|\mathbf k|_\infty\le K$. The first estimate bounds the gap
between $u_K$ and the exact solution $u$.

\begin{proposition}[Consistency of the coefficient-regularized PDE]
\label{prop:u_u_K}
Let $\mathbb T^d=(\mathbb R/2\pi\mathbb Z)^d$ and assume $a,r$ are bounded and
Lipschitz with a uniform constant $C_{b,\mathrm{lip}}$, i.e.\ for all $x,y\in\mathbb R$
\[
|a(x)|,\,|r(x)|\le C_{b,\mathrm{lip}},\qquad
|a(x)-a(y)|,\,|r(x)-r(y)|\le C_{b,\mathrm{lip}}\,|x-y|.
\]
Fix $T<\infty$ and $m>d/2+1$. Suppose that on $[0,T)$ the exact solution $u$ of the
nonlinear PDE~\eqref{equ:reaction-diffusion} with periodic boundary conditions and
the solution $u_K$ of the coefficient-regularized PDE~\eqref{equ:projected} exist
and are unique (recall that \eqref{equ:projected} is not a spectral Galerkin
truncation: $\mathcal P_K$ acts only on the coefficient arguments and $u_K$ need not
be band-limited), and that $u\in L^\infty\bigl(0,T;H^m(\mathbb T^d)\bigr)$, so that
$H^m\hookrightarrow W^{1,\infty}\hookrightarrow L^\infty$. Then the error
$e_K=u-u_K$ satisfies, for all $t\le T$,
\begin{equation}\label{eq:prop-uuK-bound}
\|e_K(t)\|_{L^2}^2\;\le\; C_T\Bigl(\|e_K(0)\|_{L^2}^2
+ K^{-2m}\,\|u\|_{L^2(0,t;H^m)}^2\Bigr),
\end{equation}
where $C_T>0$ depends only on $T$, $D$, and the coefficient constants. The proof is
given in Appendix~\ref{app:convergence}.
\end{proposition}

\paragraph{The continuous-time branching measure.}
Let $S_K(N,t)=\{\bX^N_{i,K}(t)\}_{i\in\mathcal I(N,t)}$ be the particle set at time
$t$, indexed by $\mathcal I(N,t)$, evolving under the continuous-time limit of the
algorithm: the advection--diffusion step \eqref{eq:euler_maruyama} is realized as
the SDE~\eqref{equ:overd_langevin}, and the reaction step
(Section~\ref{sec:reaction}) as a birth--death process with instantaneous rates
$c^\pm:=r^\pm(\mathcal P_K\mu^N)$, where $r^\pm=\max(\pm r,0)$. The associated
empirical (finite) measure is
\[
\mu^N(t,\intd\bx)=\frac1N\sum_{i\in\mathcal I(N,t)}\delta_{\bX^N_{i,K}(t)}(\intd\bx).
\]
By It\^o's formula, for every $f\in C^2(\mathbb T^d)$,
\begin{equation}\label{eq:weak-form}
\begin{aligned}
\langle f,\mu^N(t)\rangle
&= \langle f,\mu^N(0)\rangle
 - \int_0^t \langle\nabla f,\,a(\mathcal P_K\mu^N)\mu^N\rangle\,\intd\tau
 + \int_0^t \langle\Delta f,\,D\,\mu^N\rangle\,\intd\tau\\
&\quad + \int_0^t \langle f,\,r(\mathcal P_K\mu^N)\mu^N\rangle\,\intd\tau
 + \frac1N\int_0^t\sum_{i\in\mathcal I(N,\tau)}\sqrt{2D}\,
   \nabla f(\bX_{i,K}^N(\tau))\cdot\intd\bW^N_i(\tau)\\
&\quad + \frac1N\int_0^t\sum_{i\in\mathcal I(N,\tau)}
   f(\bX_{i,K}^N(\tau))\,\intd\widetilde{\mathcal R}^{+,N}_i(\tau)
 - \frac1N\int_0^t\sum_{i\in\mathcal I(N,\tau)}
   f(\bX_{i,K}^N(\tau))\,\intd\widetilde{\mathcal R}^{-,N}_i(\tau),
\end{aligned}
\end{equation}
where the last two terms are the compensated Poisson martingales associated with the
birth ($+$) and death ($-$) counts $\mathcal R^{\pm,N}_i$ of rate $c^\pm$. Their
net compensator equals $c^+-c^-=r$, which is exactly the deterministic reaction
drift already present in~\eqref{eq:weak-form}, so compensation records the reaction
once as drift and once as a mean-zero martingale; the precise definitions and this
bookkeeping are recalled in Appendix~\ref{app:convergence}. The next estimate
controls the gap between $u_K$ and $\mu^N$ in the negative norm $H^{-s}$,
$s>d/2+1$, with a constant independent of $N$.

\begin{proposition}[Continuous-time particle approximation]
\label{prop:u_K_mu}
Under the coefficient assumptions of Proposition~\ref{prop:u_u_K}, fix $s>d/2+1$ and
the cut-off $K$, and suppose that on $[0,T)$ the active set is uniformly bounded,
$|\mathcal I(N,t)|\le C_N N$, the projected PDE~\eqref{equ:projected} has a unique
solution $u_K$, and the band-limited coefficients and the solution are uniformly
$W^{\lceil s\rceil,\infty}$-regular,
\begin{equation}\label{eq:prop2-coeff}
\sup_{t\le T}\Bigl(\bigl\|a(\mathcal P_K\mu^N(t))\bigr\|_{W^{\lceil s\rceil,\infty}}
+\bigl\|r(\mathcal P_K\mu^N(t))\bigr\|_{W^{\lceil s\rceil,\infty}}
+\|u_K(t)\|_{W^{\lceil s\rceil,\infty}}\Bigr)\le C_{T,K}.
\end{equation}
At fixed $K$ this is automatic: $\mathcal P_K\mu^N$ is a trigonometric polynomial of
degree $\le K$, so the inverse estimate gives $\|\mathcal P_K\mu^N\|_{W^{m,\infty}}
\le C_{m,K}\|\mu^N\|_{H^{-s}}\le C_{m,K}C_N\sqrt{C_s}$, while $a,r$ are smooth
functions of $\mathcal P_K\mu^N$ and $u_K$ is the smooth solution
of~\eqref{equ:projected}. Writing $E_K(t)=\tfrac12\|u_K(t)-\mu^N(t)\|_{H^{-s}}^2$,
there is a constant $C=C_{T,K}$ depending on $C_N,C_{b,\mathrm{lip}},D,s,d$ and on
$K$ through~\eqref{eq:prop2-coeff} but \emph{independent of $N$} such that
\begin{equation}\label{eq:prop-uKmu-diff}
\frac{d}{dt}\,\mathbb E E_K(t)\le C\,\mathbb E E_K(t)+\frac{C}{N},
\end{equation}
and therefore, by Gr\"onwall with constant forcing,
\begin{equation}\label{eq:prop-uKmu-bound}
\mathbb E E_K(t)\le e^{Ct}\,\mathbb E E_K(0)+\frac{e^{Ct}-1}{N}
\le e^{Ct}\,\mathbb E E_K(0)+\frac{C_T}{N},\qquad t\le T.
\end{equation}
In particular, if $\mu^N(0)$ is an unbiased sample of $u_K(0)$ so that
$\mathbb E E_K(0)\le C_s/N$, then $\mathbb E\|u_K(t)-\mu^N(t)\|_{H^{-s}}^2\le
C_{T,K}/N$ for all $t\le T$. The proof is given in
Appendix~\ref{app:convergence}.
\end{proposition}

Combining the two propositions gives the continuous-time convergence of the
branching measure to the exact solution.

\begin{theorem}[Continuous-time convergence]
\label{thm:abstract_error}
Let the hypotheses of Propositions~\ref{prop:u_u_K} and~\ref{prop:u_K_mu} hold with
$m>d/2+1$ and $s>d/2+1$. Assume the scheme is initialized consistently and without
bias: $u_K(0)=u(0)$ and $\mu^N(0)$ is an i.i.d.\ unbiased sample of $u(0)$, so that
$\mathbb E\|u(0)-\mu^N(0)\|_{H^{-s}}^2\le C_s/N$. Then for all $t_n\le T$
\begin{equation}\label{eq:u-muN-cont}
\mathbb E\bigl\|u(t_n)-\mu^N(t_n)\bigr\|_{H^{-s}}^2
\;\le\; C_{T,K}\Bigl(K^{-2m}\,\|u\|_{L^2(0,t_n;H^m)}^2+\tfrac1N\Bigr),
\end{equation}
with $C_{T,K}$ independent of $N$. The proof is given in
Appendix~\ref{app:convergence}.
\end{theorem}

\begin{remark}[Fully discrete rate and the role of the kernel]\label{rem:rate}
The implemented algorithm returns the fully discrete measure $\mu^N_n$, which
differs from $\mu^N$ by the time-discretization term of~\eqref{eq:err-decomp}. That
term is controlled, for fixed $K$, by $\mathbb E\|\mu^N(t_n)-\mu^N_n\|_{H^{-s}}^2\le
C_{T,K}\tau^2$ in the companion paper described above; combining it with
Theorem~\ref{thm:abstract_error} through the triangle inequality gives the fully
discrete rate
\begin{equation}\label{equ:abstract_rate_full}
\Bigl(\mathbb E\bigl\|u(t_n)-\mu^N_n\bigr\|_{H^{-s}}^2\Bigr)^{1/2}
\;\le\; C_{T,K}\bigl(N^{-1/2}+\tau+K^{-m}\bigr),
\end{equation}
with Monte-Carlo ($N^{-1/2}$), time-splitting ($\tau$), and reconstruction
($K^{-m}$) contributions; the manufactured-solution study of Section~\ref{sec:mms}
probes these three terms. The continuous-time bound~\eqref{eq:u-muN-cont}, and hence
the $N^{-1/2}$ and $K^{-m}$ rates, are identical for the Poisson and
minimum-variance kernels: as $\tau\downarrow0$ both converge to the same
Poisson-clock birth--death limit, and their conditional offspring variances differ
only at order $\tau^2$. The advantage of the minimum-variance kernel is the smaller
conditional variance of Proposition~\ref{prop:stochastic_branch} and the resolution
gains documented in Section~\ref{sec:branch_vs_weight}, not a faster asymptotic
rate.
\end{remark}

\begin{remark}[Scope of the analysis]\label{rem:scope}
The constant $C_{T,K}$ depends on the horizon $T$ and grows with $K$ through the
Lipschitz and inverse constants of the band-limited coefficients, so it is not
uniform in $K$; \eqref{equ:abstract_rate_full} should be read as convergence in $N$
and $\tau$ at fixed $K$, or under a joint regime in which $K$ grows slowly enough
that $C_{T,K}K^{-m}\to0$. The estimates hold under the regularity and localization
hypotheses of Propositions~\ref{prop:u_u_K}--\ref{prop:u_K_mu}, in particular before
any singular blow-up time, and do not assert convergence through a Dirac
concentration. Finally, Theorem~\ref{thm:abstract_error} controls the $H^{-s}$
distance of the empirical measure; at fixed $K$ this is equivalent, up to a
$K$-dependent constant, to the $L^2$ error of the degree-$K$ reconstructed field of
Section~\ref{sec:interpolation} reported in Section~\ref{sec:numerics}, because all
$H^\sigma$ norms are equivalent on the degree-$K$ trigonometric class.
\end{remark}


\section{Coupled systems and implementation protocol}
\label{sec:systems}
\label{sec:extension}

The scalar construction in Section~\ref{sec:method} is used in the numerical experiments through the following finite-measure formulation.  Consider a system with components $U=(u_1,\ldots,u_L)$ that can be written, for each component, in the multiplicative reaction form
\begin{equation}
\partial_t u_\ell
= -\nabla\!\cdot\bigl(b_\ell[U]u_\ell\bigr)
  + D_\ell \Delta u_\ell
  + r_\ell[U]u_\ell,
\qquad \ell=1,\ldots,L .
\label{equ:system_multiplicative}
\end{equation}
The drift $b_\ell[U]$ may depend nonlocally on all components; for example, in Keller--Segel it is obtained from a chemical field solve, and in the kinetic test it depends on a spatial marginal.  Each component is represented by a finite empirical measure
\begin{equation}
\mu_{\ell,n}^N
= \omega_{\ell,0}\sum_{i\in \mathcal I_{\ell,n}} w_{\ell,i,n}\,\delta_{X_{\ell,i,n}},
\qquad
\omega_{\ell,0}=\frac{M_{\ell,0}}{N_{\ell,0}},
\qquad
M_{\ell,0}=\int u_\ell(0,x)\,dx .
\label{equ:system_measure}
\end{equation}
For the equal-weight branching method all active particles in component $\ell$ have the same weight $w_{\ell,i,n}\equiv 1$; mass change is recorded by the active count $|\mathcal I_{\ell,n}|$.  For the weighted baseline the active count is fixed and $w_{\ell,i,n}$ is updated multiplicatively.  In both cases the reconstructed physical field is
\begin{equation}
\widetilde u_{\ell,n}=P_K\mu_{\ell,n}^N,
\label{equ:system_reconstruction}
\end{equation}
where $P_K$ denotes the Fourier reconstruction of Section~2.5, or the lower-dimensional field reconstruction specified in the experiment.  The normalization in \eqref{equ:system_measure} is important: the reconstruction carries the current physical mass and therefore can be used directly in the drift and reaction coefficients.

One time step is a direct extension of the scalar splitting.  First the fields are reconstructed from the current empirical measures.  Then each component is transported by the SDE associated with its drift--diffusion operator.  The fields are reconstructed again at the transported positions, and the reaction multiplier $m_{\ell,i}=\exp\{r_\ell[\widetilde U^*](X^*_{\ell,i})\tau\}$ is represented either as a weight update or by one of the integer branching kernels of Section~2.4.  This is the algorithm actually used by the scalar, Keller--Segel, and kinetic experiments; the only difference between examples is how the fields entering $b_\ell$ and $r_\ell$ are reconstructed.

\begin{algorithm}[t]
\caption{One split step for a coupled finite-measure system}
\label{alg:system_step}
\begin{algorithmic}[1]
\Require particle sets $\{X_{\ell,i,n},w_{\ell,i,n}\}_{i\in\mathcal I_{\ell,n}}$ for $\ell=1,\ldots,L$; step $\tau$; reconstruction operator $P_K$; reaction representation $\mathsf{kernel}\in\{\mathsf{weighted},\mathsf{Poisson},\mathsf{minvar}\}$
\State reconstruct $\widetilde U_n=(P_K\mu_{1,n}^N,\ldots,P_K\mu_{L,n}^N)$
\For{$\ell=1,\ldots,L$}
\State transport each active particle by an Euler--Maruyama, exact OU, or problem-specific SDE step for $b_\ell[\widetilde U_n]$ and $D_\ell$:
\[
X^*_{\ell,i}=X_{\ell,i,n}+b_\ell[\widetilde U_n](X_{\ell,i,n})\tau+\sqrt{2D_\ell\tau}\,\xi_{\ell,i,n}
\]
\EndFor
\State reconstruct the transported fields $\widetilde U^*=(P_K\mu_{1}^{N,*},\ldots,P_K\mu_L^{N,*})$
\For{$\ell=1,\ldots,L$ and $i\in\mathcal I_{\ell,n}$}
\State set $m_{\ell,i}=\exp\{r_\ell[\widetilde U^*](X^*_{\ell,i})\tau\}$
\If{$\mathsf{kernel}=\mathsf{weighted}$}
\State set $w_{\ell,i,n+1}=w_{\ell,i,n}m_{\ell,i}$ and keep the transported particle
\ElsIf{$\mathsf{kernel}=\mathsf{Poisson}$}
\State draw $\nu_{\ell,i}=1+\mathrm{Poisson}(m_{\ell,i}-1)$ if $m_{\ell,i}\ge 1$, and $\nu_{\ell,i}=\mathrm{Bernoulli}(m_{\ell,i})$ otherwise
\State append $\nu_{\ell,i}$ equal-weight copies of $X^*_{\ell,i}$
\Else
\State draw $\nu_{\ell,i}=\lfloor m_{\ell,i}\rfloor+\mathrm{Bernoulli}(m_{\ell,i}-\lfloor m_{\ell,i}\rfloor)$ and append $\nu_{\ell,i}$ equal-weight copies of $X^*_{\ell,i}$
\EndIf
\EndFor
\State optionally apply the explicitly reported population-control step of Section~2.6
\end{algorithmic}
\end{algorithm}

Equation~\eqref{equ:system_multiplicative} also clarifies the scope of the coupled tests.  Branching is exact in conditional mean for sources that are absolutely continuous with respect to the current component, i.e. $F_\ell[U]=r_\ell[U]u_\ell$.  A source that creates component $\ell$ in regions where $u_\ell$ is absent, such as a cross-species injection term, should instead be represented by an injection/sampling kernel.  In the present paper the main branching evidence comes from the scalar localized-growth tests and from the kinetic field-coupled test, both of which have multiplicative reaction rates.  The conservative Keller--Segel and three-dimensional focusing examples are used to test the particle transport and field solve in concentration-dominated regimes, not to claim an additional branching advantage.

\section{Numerical experiments}
\label{sec:numerics}

We test the proposed algorithm on a sequence of numerical examples.  We first verify
the method on a smooth manufactured solution and then compare equal-weight branching
with weighted and resampled weighted particles for scalar linear reaction--diffusion
problems with localized growth.  We next consider Keller--Segel concentration dynamics:
in two dimensions we use an LDG-aligned benchmark and compare the particle diagnostics
with a direct LDG reference, while in three dimensions we test focusing behavior in radial
and nonradial configurations.  Finally, we apply the method to a six-dimensional
field-coupled kinetic problem, where only low-dimensional spatial fields are reconstructed.

Throughout this section, unless otherwise specified, relative \(L^2\) errors are computed
from the reconstructed physical field.  Let
\(
\widetilde u^{N,K}(t,\mathbf x)=P_K\mu_t^N(\mathbf x)
\)
denote the reconstructed particle solution, and let \(u_{\rm ref}\) denote the reference
solution on the evaluation grid.  On a uniform grid with cell volume \(\Delta V\), we define
\begin{equation}
E_{L^2}(t)
=
\frac{
\left(\Delta V\sum_j
\left|\widetilde u^{N,K}(t,\mathbf x_j)-u_{\rm ref}(t,\mathbf x_j)\right|^2\right)^{1/2}
}{
\left(\Delta V\sum_j
\left|u_{\rm ref}(t,\mathbf x_j)\right|^2\right)^{1/2}
}.
\label{eq:num_L2_error}
\end{equation}
For a local diagnostic region \(\Omega_{\rm loc}\), the same formula is used with both
sums restricted to grid points in \(\Omega_{\rm loc}\).

\subsection{Manufactured-solution verification}
\label{sec:mms}

We first use a smooth scalar problem for which the exact solution is known.  On $\mathbb T^2=[-\pi,\pi]^2$ we solve
\begin{equation}
\partial_t u=-\nabla\cdot(\mathbf b u)+D\Delta u+r(t,\mathbf x)u,
\label{equ:mms_pde}
\end{equation}
with prescribed positive solution
\begin{equation}
 u_{\rm ex}(t,\mathbf x)=e^{\gamma t}\left(1+a_1\cos x_1+a_2\sin x_2+a_3\cos(x_1+x_2)\right).
\label{equ:mms_exact}
\end{equation}
The reaction rate is obtained by substituting \eqref{equ:mms_exact} into \eqref{equ:mms_pde},
\begin{equation}
 r(t,\mathbf x)=\frac{\partial_t u_{\rm ex}+\nabla\cdot(\mathbf b u_{\rm ex})-D\Delta u_{\rm ex}}{u_{\rm ex}}.
\label{equ:mms_r}
\end{equation}
We use $D=0.05$, $\mathbf b=(0.5,0.3)$, $\gamma=0.3$, $(a_1,a_2,a_3)=(0.2,0.15,0.1)$, and $T=0.5$.

The particle-number refinement in Table~\ref{tab:mms_N} and Fig.~\ref{fig:mms} gives a fitted slope $-0.49$, matching the Monte-Carlo term $N^{-1/2}$ in the error decomposition.  At the largest tested $N$, the direct error against the exact solution is nearly flat as $\tau$ is refined, because the time-splitting error is below the sampling floor.  To expose that term, we use a lineage-coupled $(\tau,\tau/2)$ refinement and project the paired difference onto the deterministic splitting-bias field described in Appendix~B.  The projected branching difference scales as $\tau^{0.97}$, and the dominant Fourier modes in Table~\ref{tab:mms_modes} show the same first-order behavior.  Thus this test verifies both observable pieces of the fully discrete rate: the $N^{-1/2}$ sampling fluctuation and the first-order splitting contribution.

\begin{table}[t]
\centering
\caption{Manufactured-solution particle-number study with $\tau=2.5\times10^{-3}$ and $K=8$.  Values are relative $L^2$ errors averaged over four seeds.}
\label{tab:mms_N}
\begin{tabular}{c c c c c c c}
\toprule
$N$ & $2\times10^3$ & $4\times10^3$ & $8\times10^3$ & $1.6\times10^4$ & $3.2\times10^4$ & $6.4\times10^4$ \\
\midrule
rel. $L^2$ error & 0.325 & 0.226 & 0.167 & 0.115 & 0.083 & 0.059 \\
\bottomrule
\end{tabular}
\end{table}

\begin{figure}[t]
\centering
\includegraphics[width=0.8\linewidth]{figure/mms_convergence.pdf}
\caption{Manufactured-solution verification.  Left: relative $L^2$ error versus particle number $N$, with fitted slope $-0.49$.  Right: time-step refinement.  The total error against the exact solution is at the Monte-Carlo floor; the orange curve is the lineage-coupled branching difference projected onto the splitting-bias direction and shows first-order scaling.}
\label{fig:mms}
\end{figure}

\begin{table}[t]
\centering
\caption{Mode-resolved convergence of the branching splitting bias.  For the dominant Fourier modes of the deterministic splitting-bias field $\varphi$, the spectral mean-field coefficient and the branching estimate both converge at first order in $\tau$.}
\label{tab:mms_modes}
\begin{tabular}{c c c}
\toprule
mode & spectral slope & branching slope \\
\midrule
$\sin x_1\sin x_2$ & 0.99 & 0.91 \\
$\cos x_1\cos x_2$ & 1.00 & 0.94 \\
$\cos x_1$ & 1.00 & 0.95 \\
aggregate $(\|\varphi\|,\;P=\bar\rho\|\varphi\|)$ & 1.00 & 0.97 \\
\bottomrule
\end{tabular}
\end{table}

\subsection{Stationary localized growth: reaction representation at equal transport}
\label{sec:branch_vs_weight}
To compare particle-level realizations of the reaction operator, we consider the
scalar localized-growth problem
\begin{equation}
\partial_t u=D\Delta u+r(\mathbf x)u,
\qquad
r(\mathbf x)=\lambda G(\mathbf x)-\beta,
\qquad
G(\mathbf x)=\exp\!\left(-\frac{|\mathbf x-\mathbf x_0|_{\mathbb T}^2}{2\sigma^2}\right),
\label{equ:loc_growth}
\end{equation}
with \(D=0.05\), \(\lambda=8\), \(\beta=1\), \(\sigma=0.5\), \(\mathbf x_0=0\),
\(u_0=1+G/2\), \(T=1\), \(\tau=2\times10^{-3}\), \(K=16\), and
\(N_0=2\times10^4\).  The deterministic reference is a Fourier Strang split-step
solution on a \(256^2\) grid.  The diagnostic region is the fixed growth set
\(\Omega_{\rm grow}=\{G\ge 0.5\}\).

The final-time comparison is summarized in Table~\ref{tab:bw}.  In this table,
\[
\mathrm{nESS}
=
\frac{(\sum_i w_i)^2}{N_{\rm rep}\sum_i w_i^2},
\]
where \(N_{\rm rep}\) is the number of particles in the reported representation at
\(T\); hence \(\mathrm{nESS}=1\) for equal-weight branching.  The peak error is the absolute
peak-height error,
\[
E_{\rm peak}
=
\left|
\|\widetilde u^{N,K}(T)\|_{L^\infty}
-
\|u_{\rm ref}(T)\|_{L^\infty}
\right|.
\]  The raw weighted
method develops the expected localized weight degeneracy: its normalized ESS is
\(0.021\), the maximum-to-mean weight ratio is \(133\), and the relative \(L^2\) error is
\(0.269\pm0.041\).  ESS-triggered resampling reduces the weight degeneracy and improves
the error, but it keeps a fixed active population.  Equal-weight branching instead
represents the same reaction-induced mass growth as particle count.  At the same initial
budget, the Poisson kernel gives \(0.062\pm0.035\), while the minimum-variance kernel
gives \(0.033\pm0.011\), consistent with its smaller conditional offspring variance.
The corresponding reconstructed fields and error histories are shown in
Fig.~\ref{fig:bw_snap_l2}; the raw weighted reconstruction under-resolves the peak,
whereas ESS resampling and minimum-variance branching recover it more accurately.

For the cost comparison, we denote the active particle count by
\(
N_{\rm act}(t)=|\mathcal I(t)|
\)
and use the integrated particle-step count
\(
W_{\rm p}=\sum_{n=0}^{T/\tau}N_{\rm act}(t_n).
\)
For weighted and resampled weighted particles \(N_{\rm act}(t)\) is fixed by the chosen
particle budget, whereas for branching it is generated by the integer reaction
realization.  The cost controls in Table~\ref{tab:bw_cost} separate the representation
effect from a simple increase in particle budget.  At matched work
\(W_{\rm p}\approx1.9\times10^7\), ESS-triggered systematic resampling gives
\(0.047\pm0.012\), while minimum-variance branching gives \(0.033\pm0.011\).  The larger
\(1.3\times10^5\)-particle ESS run is more accurate, but it uses about
\(6.6\times10^7\) particle steps.  Thus the comparison should be read as a controlled
representation result, not as a claim that branching dominates every possible resampling
strategy at every budget.

\begin{table}[t]
\centering
\caption{Stationary localized-growth benchmark at \(T=1\), with common transport and
\(N_0=2.0\times10^4\).  Weighted rows report normalized ESS and max:mean weight ratio.
Branching rows are equal-weight, so nESS is one.  The capped row uses minimum-variance
branching with \(N_{\max}=4.0\times10^4\).}
\label{tab:bw}
\begin{tabular}{l c c c c}
\toprule
method & rel. \(L^2\) error & nESS & max:mean \(w\) & \(E_{\rm peak}\) \\
\midrule
weighted & \(0.269\pm0.041\) & 0.021 & 133 & 149 \\
weighted + ESS & \(0.073\pm0.033\) & 0.56 & 3.1 & \(37\pm34\) \\
Poisson br. & \(0.062\pm0.035\) & 1 & -- & 31 \\
min.-var. br. & \(\mathbf{0.033\pm0.011}\) & 1 & -- & \(\mathbf{12.7}\) \\
capped min.-var. & \(0.048\pm0.010\) & 1 & -- & 17.4 \\
\bottomrule
\end{tabular}
\end{table}

\begin{table}[t]
\centering
\caption{Cost and large-budget controls for the stationary localized-growth benchmark.
Here \(W_{\rm p}=\sum_n N_{\rm act}(t_n)\) is the integrated particle-step count.}
\label{tab:bw_cost}
\begin{tabular}{l c c c}
\toprule
method & initial budget & \(W_{\rm p}\) & rel. \(L^2\) error \\
\midrule
weighted & \(6.0\times10^4\) & \(3.0\times10^7\) & \(0.151\pm0.021\) \\
weighted & \(1.3\times10^5\) & \(6.6\times10^7\) & \(0.089\pm0.018\) \\
weighted + ESS & \(3.8\times10^4\) & \(1.9\times10^7\) & \(0.047\pm0.012\) \\
weighted + ESS & \(1.3\times10^5\) & \(6.6\times10^7\) & \(0.028\pm0.006\) \\
min.-var. br. & \(2.0\times10^4\) & \(\mathbf{1.9\times10^7}\) & \(\mathbf{0.033\pm0.011}\) \\
\bottomrule
\end{tabular}
\end{table}

\begin{figure}[t]
\centering
\includegraphics[width=0.4\linewidth]{figure/snapshots_final.pdf}\quad
\includegraphics[width=0.4\linewidth]{figure/l2_vs_t.pdf}
\caption{Stationary localized-growth benchmark.  Left: final-time reconstructed fields
at \(T=1\) for the deterministic reference, raw weighted particles, ESS-resampled
weighted particles, and minimum-variance branching, all with \(N_0=2.0\times10^4\).
Right: relative \(L^2\) error versus time for the same initial budget.  Raw weighting
under-resolves the localized peak, while ESS resampling and minimum-variance branching
reduce the error; quantitative final-time diagnostics are reported in
Table~\ref{tab:bw}.}
\label{fig:bw_snap_l2}
\end{figure}


\subsection{Switching growth under a moving source}
\label{sec:switch}
We further test the representation comparison with a growth region that switches from $\mathbf x_A=(-1.2,0)$ to $\mathbf x_B=(1.2,0)$ at mid-time:
\begin{equation}
\partial_tu=D\Delta u+r(t,\mathbf x)u,
\qquad
r(t,\mathbf x)=\lambda\bigl[(1-s(t))G_\sigma(\mathbf x;\mathbf x_A)+s(t)G_\sigma(\mathbf x;\mathbf x_B)\bigr]-\beta,
\qquad
s(t)=\frac12\left(1+\tanh\frac{t-T/2}{\delta}\right).
\label{equ:switch}
\end{equation}
Here $u_0\equiv1$, $D=0.02$, $\lambda=10$, $\beta=0.1$, $\sigma=0.25$, $T=1.2$, $\delta=0.04$, $\tau=10^{-3}$, $K=48$, and $N_0=2\times10^4$.  The reference field \(u_{\rm ref}\) is computed by a deterministic Fourier pseudospectral solver on a \(512\times512\) collocation grid.  
The second-stage region $B_B$ is the discriminating diagnostic.  As shown in Table~\ref{tab:switch}, both branching kernels outperform the tested global resampling baselines globally and in $B_B$.  ESS-triggered resampling is best only in the first-stage region $B_A$, where it has concentrated the fixed particle budget before the switch.  By $t=T/2^-$, global resampling has reduced the number of distinct surviving ancestors to $63$--$68\%$ of $N_0$, while branching retains about $95\%$ and has grown the active population to $4.2\times10^4$.  Thus the failure mode in this test is not only weight variance; it is loss of future-growth lineages.

\begin{table}[t]
\centering
\caption{Switching-growth benchmark at $T=1.2$ (eight seeds): relative $L^2$ errors globally and in the first- and second-stage regions, active particle count, and fraction of distinct surviving ancestors at the switch.}
\label{tab:switch}
\begin{tabular}{l c c c c c}
\toprule
method & global $L^2$ & $B_A$ $L^2$ & $B_B$ $L^2$ & $N_{\rm act}$ & distinct anc. at switch \\
\midrule
weighted & $0.900\pm0.048$ & $1.014\pm0.120$ & $0.690\pm0.039$ & $2.0\times10^4$ & 100\% \\
weighted + resample (ESS) & $0.288\pm0.025$ & $\mathbf{0.162\pm0.027}$ & $0.297\pm0.038$ & $2.0\times10^4$ & 68\% \\
weighted + resample (every step) & $0.271\pm0.062$ & $0.234\pm0.087$ & $0.226\pm0.107$ & $2.0\times10^4$ & 63\% \\
Poisson branching & $\mathbf{0.199\pm0.045}$ & $0.183\pm0.063$ & $\mathbf{0.168\pm0.067}$ & $4.2\times10^4$ & 95\% \\
minimum-variance branching & $0.203\pm0.051$ & $0.170\pm0.136$ & $0.174\pm0.022$ & $4.2\times10^4$ & 95\% \\
\bottomrule
\end{tabular}
\end{table}

\begin{figure}[t]
\centering
\includegraphics[width=0.8paper\linewidth]{figure/switch_snapshots_magnifier_4x1.png}
\caption{Switching-growth benchmark at $T=1.2$, after growth has moved from $\mathbf x_A$ to $\mathbf x_B$.  The weighted reconstruction is dominated by a few large weights; systematic resampling and minimum-variance branching both recover smoother fields, with the smallest second-stage error obtained by branching.}
\label{fig:switch}
\end{figure}

\subsection{Two-dimensional Keller--Segel: coupled solver and pre-singular concentration}
\label{sec:ks}

We next move to Keller--Segel dynamics in physical dimension two.  These experiments test the particle-field coupling and the behavior of the mesh-free representation in concentration-dominated regimes.  They are deliberately interpreted as pre-singular concentration tests, not as convergence through blow-up, because the Fourier reconstruction is bandwidth-limited while the continuum Keller--Segel dynamics may approach singular measures in two dimensions~\cite{herrero1997blow,perthame2007transport,giga2011asymptotic,souplet2019blow,filbet2006finite,chertock2008second,chertock2018high,shen2020unconditionally,chen2022error,liu2018positivity,zhong2024direct,li2017local,epshteyn2009new,epshteyn2019efficient,strehl2013positivity}.

\paragraph{Coupled parabolic--parabolic implementation check.}
We first solve
\begin{equation}
 u_t-\nabla\cdot(\nabla u-u\nabla v)=0,
 \qquad
 v_t-\Delta v=u-v,
 \qquad \mathbf x\in[0,2\pi]^2,
\label{equ:case1}
\end{equation}
with periodic boundary conditions and initial data $u_0(\mathbf x)=\sin^2(x_1)\cos^2(x_2)$ and $v_0(\mathbf x)=\cos x_1+\cos x_2+2$.  Since the cell equation is conservative, $\int u(t,x)dx=\pi^2$, and the chemical mass satisfies
\[
M_v'(t)=-M_v(t)+\pi^2,
\qquad
M_v(t)=\pi^2+[M_v(0)-\pi^2]e^{-t}.
\]
Figure~\ref{fig:total_mass} shows that the computed chemical mass follows this identity to sampling accuracy.  Figures~\ref{fig:fd_comparison_u}--\ref{fig:fd_comparison_v} compare the particle fields with a finite-difference reference, and Fig.~\ref{fig:convergence_error} shows the expected $N^{-1/2}$ trend against a high-particle reference.  This is an implementation check for the coupled transport and field reconstruction, not the main branching comparison.

\begin{figure}[h]
\centering
\includegraphics[width=0.8\linewidth]{figure/total_mass.pdf}
\caption{Mass-balance test for the conservative Keller--Segel system.  The chemical mass follows the exact balance law to sampling accuracy.}
\label{fig:total_mass}
\end{figure}

\begin{figure}[h]
\centering
\includegraphics[width=\linewidth]{figure/case1_u.pdf}
\caption{Conservative Keller--Segel test: comparison of the cell field $u$ from the particle method and from a finite-difference reference.}
\label{fig:fd_comparison_u}
\end{figure}

\begin{figure}[h]
\centering
\includegraphics[width=\linewidth]{figure/case1_v.pdf}
\caption{Conservative Keller--Segel test: comparison of the chemical field $v$ from the particle method and from a finite-difference reference.}
\label{fig:fd_comparison_v}
\end{figure}

\begin{figure}[h]
\centering
\includegraphics[width=0.8\linewidth]{figure/case1_err.pdf}
\caption{Relative $L^2$ error of $u$ and $v$ versus initial particle count.  The reference solution uses $N_0=320{,}000$, and both fields show the expected sampling trend.}
\label{fig:convergence_error}
\end{figure}

\paragraph{Supercritical concentration.}
We then use the concentrated Gaussian data
\[
 u_0(\mathbf x)=840\exp[-84(x_1^2+x_2^2)],
 \qquad
 v_0(\mathbf x)=420\exp[-42(x_1^2+x_2^2)],
\]
whose cell mass is $10\pi$, above the two-dimensional critical mass $8\pi$.  This profile is a standard test in grid-based Keller--Segel studies, including LDG and finite-volume methods~\cite{li2017local,epshteyn2009new,chertock2008second}.  We use it to ask whether the particle cloud tracks the concentrating region without an externally imposed adaptive mesh.  Figures~\ref{fig:case2_u} and~\ref{fig:comparison} show rapid peak formation and compare the particle reconstruction with uniform-grid finite-difference solutions.  The finite-difference peak sharpens with mesh refinement, while the particle reconstruction is stable over the tested particle counts.  This comparison is qualitative; specialized positivity-preserving adaptive grid methods remain the reference for resolving the singular limit.

Rather than fit a blow-up time, we report grid-free concentration diagnostics in Fig.~\ref{fig:ks_concentration_diag}.  The second moment about the cell centroid decreases by about a factor of six, and the local particle count in a fixed ball of radius $0.05$ grows by about $4.5$ times.  The curves overlap for $N=4\times10^4$, $1.6\times10^5$, and $6.4\times10^5$, indicating that the observed concentration is not a sampling artifact.

\begin{figure}
\centering
\includegraphics[width=0.8\linewidth]{figure/case2_u.pdf}
\caption{Pre-singular concentration test: particle approximation of $u$ at early times.  The cell density concentrates rapidly while remaining stable over the tested particle counts.}
\label{fig:case2_u}
\end{figure}

\begin{figure}
\centering
\includegraphics[width=0.85\linewidth]{figure/comparison.pdf}
\caption{Pre-singular concentration test at $t=5\times10^{-5}$: particle approximation compared with finite-difference approximations on three uniform grids.  The particle peak is stable across particle counts; the uniform-grid solution sharpens with mesh refinement.}
\label{fig:comparison}
\end{figure}

\begin{figure}[t]
\centering
\includegraphics[width=\linewidth]{figure/ks_concentration_diag.pdf}
\caption{Grid-free concentration diagnostics.  Left: the second moment of the cell cloud decreases by about a factor of six.  Right: the local particle count in a fixed ball of radius $0.05$ about the centroid grows by about $4.5\times$.  These diagnostics measure approach toward concentration, not a blow-up time.}
\label{fig:ks_concentration_diag}
\end{figure}

\paragraph{Core-local resolution diagnostics.}
For a more resolution-sensitive diagnostic, we also run the parabolic--elliptic reduction $-\Delta v+v=u$ on a core-adaptive Fourier window, using the same concentrated initial profile.  We track the reconstructed peak $P_{K,N}(t)=\|P_K\mu_t^N\|_\infty$, the reconstructed norm $S_{K,N}(t)=\|P_K\mu_t^N\|_{L^2}$, and reconstruction-free core radii $R_q(t)$.  Snapshots are recorded at $t=5\times10^{-5}$, $10^{-4}$, $1.5\times10^{-4}$, $5\times10^{-4}$, $10^{-3}$, $2\times10^{-3}$, and $4\times10^{-3}$.

Figures~\ref{fig:ldg_snapshots}--\ref{fig:ldg_curves} show that the inner half-mass radius decreases from about $0.091$ to about $0.06$ by $t=4\times10^{-3}$.  The reconstructed peak grows with bandwidth, approximately $1.7\times10^3$ for $K=5$ and $2.8\times10^3$ for $K=10$, so the peak itself is not treated as a converged blow-up diagnostic.  To quantify the onset of under-resolution we report
\[
 t_{\rm gap}(N,K)=\inf\{t:S_{2K,4N}(t)/S_{K,N}(t)\ge\theta\}.
\]
Table~\ref{tab:ldg_tgap} should be read exactly as this definition states: it marks when the reconstructed $L^2$ norm becomes sensitive to resolution, not when the PDE blows up.

\begin{figure}[t]
\centering
\includegraphics[width=\linewidth]{figure/ldg_snapshots.pdf}
\caption{Core-adaptive parabolic--elliptic Keller--Segel diagnostics: reconstructed contours of $u$ over the reporting times.  The field sharpens into a concentrating peak, consistent with the grid-based concentration literature.}
\label{fig:ldg_snapshots}
\end{figure}

\begin{figure}[t]
\centering
\includegraphics[width=\linewidth]{figure/ldg_curves.pdf}
\caption{Resolution diagnostics for the refined $(N,K)=(6.4\times10^5,10)$ run: reconstructed $L^2$ norm, reconstructed peak, and squared core radii.  The core radius collapse is reconstruction-free; the peak remains bandwidth-sensitive.}
\label{fig:ldg_curves}
\end{figure}

\begin{table}[t]
\centering
\caption{Resolution-gap focusing indicator $t_{\rm gap}(N,K)=\inf\{t:S_{2K,4N}(t)/S_{K,N}(t)\ge\theta\}$ for two refinement pairs.  The indicator marks the onset of under-resolution, not a blow-up time.}
\label{tab:ldg_tgap}
\begin{tabular}{c c c c c}
\toprule
base $(N,K)$ & refined $(4N,2K)$ & $t_{\rm gap}(\theta=1.05)$ & $t_{\rm gap}(\theta=1.10)$ & $\max_t S_{2K,4N}/S_{K,N}$ \\
\midrule
$(4\times10^4,5)$ & $(1.6\times10^5,10)$ & $3.0\times10^{-3}$ & -- & 1.08 \\
$(1.6\times10^5,5)$ & $(6.4\times10^5,10)$ & $2.5\times10^{-3}$ & $3.8\times10^{-3}$ & 1.11 \\
\bottomrule
\end{tabular}
\end{table}

\subsection{Three-dimensional Keller--Segel focusing}
\label{sec:ks3d}

The three-dimensional Keller--Segel experiment is a conservative particle-field focusing test.  It exercises the mesh-free transport and the three-dimensional screened-Poisson field solve, but it is not used to claim a universal critical mass or a singular blow-up time.  On $[-L/2,L/2]^3$ with $L=12$, we solve
\begin{equation}
\partial_t\rho=\Delta\rho-\chi\nabla\cdot(\rho\nabla c),
\qquad
-\Delta c+\kappa^2c=\rho-\bar\rho,
\qquad
\chi=1,
\qquad
\kappa=0.1 .
\label{equ:ks3d}
\end{equation}
The particle dynamics are $dX=\chi\nabla c\,dt+\sqrt{2}\,dW$, and $c$ is reconstructed by a three-dimensional spectral screened-Poisson solve with zero mean.  Unless otherwise stated, $\tau=5\times10^{-4}$, $T=1$, and the field bandwidth is $H=16$.

We first sweep the radial Gaussian family
\[
 \rho_0^M(\mathbf x)=M(2\pi\sigma^2)^{-3/2}\exp[-|\mathbf x|^2/(2\sigma^2)],
 \qquad
 \sigma=0.45,
 \qquad
 M\in\{20,40,60,80,100\}.
\]
The low-mass cases disperse, while the higher-mass cases focus.  The final half-mass radius is $2.07$ for $M=20$ and $1.71$ for $M=40$, but about $0.103$, $0.084$, and $0.074$ for $M=60$, $80$, and $100$, respectively (Fig.~\ref{fig:ks3d_core_radii}).  Thus the transition in this family lies between $M=40$ and $M=60$.  The core-density surrogate in Fig.~\ref{fig:ks3d_rho_core} grows rapidly for the focusing cases, but its absolute value is bandwidth- and time-step-sensitive and should be read only as a focusing indicator.

At $M=80$, the half-mass radius is stable under particle-number refinement: $N=10^5$, $2\times10^5$, and $4\times10^5$ all give $R_{0.5}\approx0.084$ at final time (Fig.~\ref{fig:ks3d_selfconv}).  The reconstructed peak still grows with bandwidth, so the radius is the reliable diagnostic here.  A nonradial four-cluster datum at the vertices of a tetrahedron gives the expected mutual attraction: the minimum inter-cluster distance decreases from $2.83$ to $2.12$ over $[0,1]$ (Fig.~\ref{fig:ks3d_tetra}).

\begin{figure}[t]
\centering
\includegraphics[width=\linewidth]{figure/ks3d_core_radii.pdf}
\caption{Three-dimensional Keller--Segel focusing.  Core radii $R_{0.5}$ and $R_{0.9}$ for the radial mass sweep.  Low masses disperse, while higher masses focus, with a transition between $M=40$ and $M=60$ for this family.}
\label{fig:ks3d_core_radii}
\end{figure}

\begin{figure}[t]
\centering
\includegraphics[width=0.8\linewidth]{figure/ks3d_rho_core.pdf}
\caption{Core-density surrogate for the three-dimensional mass sweep.  Absolute peak values are bandwidth- and time-step-limited and are used only as qualitative focusing indicators.}
\label{fig:ks3d_rho_core}
\end{figure}

\begin{figure}[t]
\centering
\includegraphics[width=\linewidth]{figure/ks3d_selfconv.pdf}
\caption{Self-convergence at $M=80$.  The core radius is stable under particle-number refinement, while the reconstructed peak still grows with bandwidth $H$.}
\label{fig:ks3d_selfconv}
\end{figure}

\begin{figure}[t]
\centering
\includegraphics[width=0.8\linewidth]{figure/ks3d_tetra.pdf}
\caption{Nonradial four-cluster datum at the vertices of a tetrahedron.  The cluster centroids mutually attract, and the minimum inter-cluster distance decreases from $2.83$ to $2.12$.}
\label{fig:ks3d_tetra}
\end{figure}

\subsection{Six-dimensional field-coupled kinetic Keller--Segel}
\label{sec:kinetic_ks}

After the two- and three-dimensional Keller--Segel tests, we turn to a phase-space problem.  The purpose is not to assemble or converge a dense six-dimensional density.  Instead, the test shows that the reaction step is unchanged in six-dimensional phase space when the nonlocal field depends only on the three-dimensional spatial marginal.  Let $(\mathbf x,\mathbf v)\in\mathbb T^3\times\mathbb R^3$ and $\rho(t,\mathbf x)=\int f(t,\mathbf x,\mathbf v)\,d\mathbf v$.  The chemical field solves
\begin{equation}
 -\Delta_{\mathbf x}c+\kappa^2c=\rho-\bar\rho,
 \qquad
 \bar\rho=|\mathbb T^3|^{-1}\int_{\mathbb T^3}\rho\,d\mathbf x,
\label{equ:kinetic_ks_poisson}
\end{equation}
with zero mean, and the kinetic equation is
\begin{equation}
\partial_t f+\mathbf v\cdot\nabla_{\mathbf x}f
=\gamma_v\nabla_{\mathbf v}\cdot\bigl[(\mathbf v-\chi\nabla_{\mathbf x}c)f\bigr]
+D_v\Delta_{\mathbf v}f
+r[\rho,c](\mathbf x)f .
\label{equ:kinetic_ks}
\end{equation}
The associated particle dynamics are
\begin{equation}
 d\mathbf X=\mathbf V\,dt,
 \qquad
 d\mathbf V=-\gamma_v(\mathbf V-\chi\nabla_{\mathbf x}c)\,dt+\sqrt{2D_v}\,d\mathbf W,
\label{equ:kinetic_ks_sde}
\end{equation}
with an exact Ornstein--Uhlenbeck velocity update over one step.  The reaction rate is
\begin{equation}
 r[\rho,c]=\lambda_g S_c(c)-\alpha_\rho S_\rho(\rho)-\beta,
 \qquad
 S_c(c)=\frac12\left(1+\tanh\frac{c-c_0}{\delta_c}\right),
 \qquad
 S_\rho(\rho)=\frac{\rho_+}{\rho_+ + \rho_0}.
\label{equ:kinetic_ks_rate}
\end{equation}
Only $\rho(\mathbf x)$ and $c(\mathbf x)$ are reconstructed, using spatial bandwidth $K_x=5$; no six-dimensional tensor grid is used.

The parameters are $\gamma_v=3$, $D_v=0.6$, $\chi=3$, $\kappa=0.5$, $\lambda_g=4$, $\alpha_\rho=3$, $\beta=0.4$, $c_0=0.05$, $\delta_c=0.05$, $\rho_0=0.2$, $T=2$, $\tau=2\times10^{-3}$, and $N_0=1.2\times10^4$.  The initial distribution is a single Gaussian blob in space with Maxwellian velocities at temperature $D_v/\gamma_v=0.2$.

The cloud focuses and grows in a bounded manner.  The half-mass radius decreases from $0.768$ to about $0.575$, the reconstructed $\|\rho\|_\infty$ increases from $0.46$ to about $11.8$, and the total mass grows by roughly a factor of eight (Fig.~\ref{fig:kinetic_ks_field}).  The reaction is genuinely field-coupled: the rate is positively correlated with $c$ and $\rho$ during activation and becomes anti-correlated at late times when the crowding term dominates (Fig.~\ref{fig:kinetic_ks_coupling}).  Global observables are similar across reaction representations, so this is not a global-accuracy benchmark.  The discriminating quantity is local resolution: the weighted method ends with nESS $0.66$ and max-to-mean weight $6.5$, whereas branching keeps equal weights and grows the active population to about $9.5\times10^4$.  In the concentration region, weighted particles provide a local effective count of about $3.3\times10^3$, while branching supplies about $2.0\times10^4$ equal-weight particles (Table~\ref{tab:kinetic_ks} and Fig.~\ref{fig:kinetic_ks_repr}).

\begin{table}[t]
\centering
\caption{Six-dimensional field-coupled kinetic Keller--Segel test at $T=2$ (two-seed representation demonstration).  For the weighted method, nESS and max:mean weight measure degeneracy and local $N_B$ is local effective sample size; for branching, all weights are equal and local $N_B$ is local particle count.}
\label{tab:kinetic_ks}
\begin{tabular}{l c c c c}
\toprule
method & final nESS & max:mean $w$ & $N_{\rm act}$ & local $N_B$ \\
\midrule
weighted & 0.66 & 6.5 & $1.2\times10^4$ & $3.3\times10^3$ \\
equal-weight branching & 1 (equal) & -- & $9.5\times10^4$ & $2.0\times10^4$ \\
\bottomrule
\end{tabular}
\end{table}

\begin{figure}[t]
\centering
\includegraphics[width=\linewidth]{figure/kinetic_ks_field.pdf}
\caption{Six-dimensional field-coupled kinetic Keller--Segel test.  Total mass, $\|c\|_\infty$, $\|\rho\|_\infty$, and spatial-marginal core radii versus time.  The cloud concentrates while the mass grows in a bounded, self-limited manner.}
\label{fig:kinetic_ks_field}
\end{figure}

\begin{figure}[t]
\centering
\includegraphics[width=\linewidth]{figure/kinetic_ks_repr.pdf}
\caption{Representation diagnostics for the kinetic test.  Weighted particles develop moderate degeneracy; branching keeps nESS $=1$ and supplies roughly sixfold more local resolution in the concentration region.}
\label{fig:kinetic_ks_repr}
\end{figure}

\begin{figure}[t]
\centering
\includegraphics[width=\linewidth]{figure/kinetic_ks_coupling.pdf}
\caption{Field coupling of the reaction in the kinetic test: correlations of the per-particle rate with $c$ and $\rho$, and the mean activation $\langle S_c\rangle$.  The sign change of the correlations reflects the transition from chemical activation to crowding-limited growth.}
\label{fig:kinetic_ks_coupling}
\end{figure}

The dense four- and six-dimensional reconstruction experiment is therefore moved out of the main numerical story and kept as Appendix~\ref{app:highdim}.  It shows that dense Fourier reconstruction has the expected $N^{-1/2}$ sampling slope at fixed bandwidth, but that its constant grows rapidly with the number of modes.  This is precisely why the kinetic experiment reconstructs only low-dimensional fields and why high-dimensional extensions should use marginal, adaptive, or low-rank reconstructions rather than dense tensor-product grids.
\section{Discussion and conclusions}
\label{sec:discussion}
We have developed an equal-weight stochastic branching particle method for non-conservative
transport--reaction--diffusion equations. The evolution is split into an SDE transport step and a particle-level
reaction step, and the central observation is that the reaction term can be represented in two qualitatively
different ways: as a growing scalar weight, or as a growing integer particle count. The weighted update is exact
in expectation but lets the local effective sample size collapse where mass is created; equal-weight branching
instead converts mass creation into particle creation, so the local resolution increases exactly where it is
needed. We placed the weighted update and two branching kernels -- a Poisson birth--death kernel and a
minimum-variance integer kernel -- in a common framework, showed that all three are unbiased for the reaction
multiplier while differing in conditional variance (Proposition~\ref{prop:stochastic_branch}), and gave a short
error decomposition into Monte-Carlo, time-splitting, and reconstruction contributions
(Theorem~\ref{thm:abstract_error} proves the continuous-time Monte-Carlo and reconstruction rates, with proofs
in~\ref{app:convergence}; the time-splitting term is treated in a companion paper).

The numerical study supports a measured set of conclusions, in the order of the experiments.  First, the
manufactured-solution test confirms the Monte-Carlo $N^{-1/2}$ rate cleanly and exposes the first-order
time-splitting term through a lineage-coupled refinement; on this smooth benchmark the reconstruction term lies
below the Monte-Carlo floor.  Second, the controlled stationary localized-growth comparison isolates the effect of
the reaction representation and shows, under shared transport and randomness, that branching maintains the local
effective sample size and local particle count where the raw weighted scheme degenerates, that adding global
particles does not cure the local degeneracy, and that branching is more accurate than the tested global
systematic-resampling baseline at equal integrated particle-step work, with no resampling schedule to tune.
Third, the switching-growth test exhibits a structural failure mode of global resampling: collapsing lineage
diversity onto the first-stage region leaves it unable to rebuild the second-stage region, so the discriminating
quantity is lineage diversity rather than weight variance alone.  Fourth, a six-dimensional, field-coupled
kinetic Keller--Segel model -- nonlinear and nonlocal through a self-consistent chemical field, yet requiring only
a three-dimensional field solve -- shows that the branching step transfers unchanged to high-dimensional phase
space without assembling a dense grid; this is a nonlinear high-dimensional representation test, not a full
convergence study of the six-dimensional density.  Finally, Keller--Segel serves as a concentration-dominated
stress test in which the mesh-free representation tracks the pre-blow-up concentration without externally imposed
mesh refinement, in two dimensions and in a self-contained three-dimensional focusing family; these are focusing
and concentration stress tests, and we make no singular-blow-up-time or universal-critical-mass claim.

Several limitations frame the scope of these conclusions and point to future work. The analysis is carried out
for the projected, regularized dynamics under regularity and localization assumptions and before the singular
blow-up time; we therefore do not claim to capture the singular profile itself, and a sharp analysis of the
coupled multi-species system remains open. The reconstruction is the only step that uses smoothness and, in its
present dense Fourier form, scales as $K^d$; scaling to genuinely high dimensions calls for low-rank or adaptive
reconstructions~\cite{TangYing2025FHT,tang2025solving,Peng_Yang_tensor_arxiv_2024,novikov2021tensor,sheng2025approximation},
which we deliberately leave to future work rather than claim here. Our weighted baselines include systematic
resampling (ESS-triggered and every-step) on one benchmark; broader resampling variants (residual, stratified,
local splitting/merging) and other growth regimes remain to be compared. Extending the particle dynamics to general
boundary conditions (through stochastic killing, reflection, or mixed rules); a cross-species injection mechanism
for sources not proportional to the local field (such as the chemical production in Keller--Segel); and exploring
score-based or generative corrections, full radiative-transfer and Vlasov--Poisson kinetic models, and a rigorous
blow-up theory, are natural further directions. The contribution of the present paper is narrower and, we believe,
useful on its own: that branching is a controllable, equal-weight numerical representation of non-conservative
mass creation, and a practical alternative to weighted particles when resolution must follow the growth.

% The framework presented here offers a fresh particle-based approach to solving reaction--diffusion--advection equations. Unlike traditional particle or Lagrangian schemes, our method is not restricted to conservation laws and is able to capture aggregation and focusing phenomena in the density variable with high fidelity. This highlights its potential as a useful tool for studying a broad class of nonlinear transport and interaction problems.

% From the theoretical perspective, the current work provides heuristic intuition primarily in the linear case. Developing a rigorous convergence theory in more general nonlinear settings remains a challenging but exciting direction for future research. On the numerical side, while the interpolation of measures poses difficulties in high-dimensional problems, this limitation also points to promising avenues for methodological innovation. In particular, Fourier interpolation combined with low-rank approximations~\cite{tang2025solving} could provide a powerful mechanism for scaling the algorithm efficiently to higher dimensions.


\section*{Acknowledgments}
The work is supported in part by the National Science Foundation under Grant DMS-2143739 and the ACCESS program through allocation MTH210005.
The authors also acknowledge the support from the Institute for Cyber-Enabled Research at Michigan State University.

\section*{Declaration of competing interest}
The authors declare that they have no known competing financial interests or personal relationships that could
have appeared to influence the work reported in this paper.

\section*{Data availability}
The configuration files and scripts that reproduce every figure and table in this paper will be made publicly
available upon publication.

\appendix

\section{Proofs of the continuous-time error estimates}
\label{app:convergence}
This appendix proves Propositions~\ref{prop:u_u_K} and~\ref{prop:u_K_mu} and
Theorem~\ref{thm:abstract_error}. Throughout, $\langle\cdot,\cdot\rangle$ is the
$L^2(\mathbb T^d)$ pairing, $\|\cdot\|_{H^\sigma}$ the periodic Sobolev norm with
$(I-\Delta)^{\sigma/2}$ as Fourier multiplier, and all boundary terms in
integration by parts vanish by periodicity.

\begin{proof}[Proof of Proposition~\ref{prop:u_u_K}]
Let $e_K=u-u_K$ and $\eta(t)=\tfrac12\|e_K(t)\|_{L^2}^2$. Subtracting
\eqref{equ:projected} from \eqref{equ:reaction-diffusion} and testing against $e_K$,
\[
\eta'(t)=\underbrace{-\!\int_{\mathbb T^d}\!\nabla e_K\!\cdot\!\bigl(a(u)u-a(\mathcal P_K u_K)u_K\bigr)\intd\bx}_{E_1}
+\underbrace{D\!\int_{\mathbb T^d}\! e_K\,\Delta e_K\intd\bx}_{E_2}
+\underbrace{\int_{\mathbb T^d}\! e_K\bigl(r(u)u-r(\mathcal P_K u_K)u_K\bigr)\intd\bx}_{E_3}.
\]
By $m>d/2+1$ and $H^m\hookrightarrow L^\infty$, set
$C_u':=\sup_{t\le T}\|u(t)\|_{L^\infty}<\infty$. By the spectral accuracy of
$\mathcal P_K$ and the triangle inequality,
\begin{equation}\label{eq:app-uuK-proj}
\|u-\mathcal P_K u_K\|_{L^2}
\le\|(I-\mathcal P_K)u\|_{L^2}+\|\mathcal P_K(u-u_K)\|_{L^2}
\le C\,K^{-m}\|u\|_{H^m}+\|e_K\|_{L^2}.
\end{equation}

\emph{Diffusion.} Integration by parts gives
$E_2=-D\|\nabla e_K\|_{L^2}^2$.

\emph{Advection.} Splitting
$a(u)u-a(\mathcal P_K u_K)u_K=(a(u)-a(\mathcal P_K u_K))u+a(\mathcal P_K u_K)(u-u_K)$,
\[
\begin{aligned}
E_1
&=-\!\int\!\nabla e_K\!\cdot\!\bigl(a(u)-a(\mathcal P_K u_K)\bigr)u
  -\!\int\!\nabla e_K\!\cdot\! a(\mathcal P_K u_K)\,e_K
=:E_{11}+E_{12}.
\end{aligned}
\]
For $E_{11}$, using $|a(u)-a(\mathcal P_K u_K)|\le C_{b,\mathrm{lip}}|u-\mathcal P_K u_K|$,
the bound $\|u\|_{L^\infty}\le C_u'$, Cauchy--Schwarz, and Young's inequality with
$\epsilon>0$,
\[
E_{11}\le \|\nabla e_K\|_{L^2}\,C_{b,\mathrm{lip}}C_u'\,\|u-\mathcal P_K u_K\|_{L^2}
\le \tfrac{\epsilon}{2}C_{b,\mathrm{lip}}C_u'\|\nabla e_K\|_{L^2}^2
+\tfrac{1}{2\epsilon}C_{b,\mathrm{lip}}C_u'\|u-\mathcal P_K u_K\|_{L^2}^2 .
\]
For $E_{12}$, using $\|a(\mathcal P_K u_K)\|_{L^\infty}\le C_{b,\mathrm{lip}}$,
\[
E_{12}\le C_{b,\mathrm{lip}}\|\nabla e_K\|_{L^2}\|e_K\|_{L^2}
\le \tfrac{\epsilon}{2}C_{b,\mathrm{lip}}\|\nabla e_K\|_{L^2}^2
+\tfrac{1}{2\epsilon}C_{b,\mathrm{lip}}\|e_K\|_{L^2}^2 .
\]
Inserting \eqref{eq:app-uuK-proj} and $(p+q)^2\le2p^2+2q^2$ into $E_{11}$,
\[
E_1\le \tfrac{\epsilon}{2}C_{b,\mathrm{lip}}(C_u'+1)\,\|\nabla e_K\|_{L^2}^2
+C_1\bigl(\|e_K\|_{L^2}^2+K^{-2m}\|u\|_{H^m}^2\bigr),
\]
where $C_1$ depends only on $C_{b,\mathrm{lip}},C_u',\epsilon$.

\emph{Reaction.} Splitting
$r(u)u-r(\mathcal P_K u_K)u_K=(r(u)-r(\mathcal P_K u_K))u+r(\mathcal P_K u_K)(u-u_K)$
and using $|r(u)-r(\mathcal P_K u_K)|\le C_{b,\mathrm{lip}}|u-\mathcal P_K u_K|$,
$\|r(\mathcal P_K u_K)\|_{L^\infty}\le C_{b,\mathrm{lip}}$, Young's inequality, and
\eqref{eq:app-uuK-proj},
\[
E_3\le C_{b,\mathrm{lip}}C_u'\|e_K\|_{L^2}\|u-\mathcal P_K u_K\|_{L^2}
+C_{b,\mathrm{lip}}\|e_K\|_{L^2}^2
\le C_3\bigl(\|e_K\|_{L^2}^2+K^{-2m}\|u\|_{H^m}^2\bigr),
\]
with $C_3$ depending only on $C_{b,\mathrm{lip}},C_u'$.

\emph{Conclusion.} Summing and using $E_2=-D\|\nabla e_K\|_{L^2}^2$,
\[
\eta'(t)\le\Bigl(\tfrac{\epsilon}{2}C_{b,\mathrm{lip}}(C_u'+1)-D\Bigr)\|\nabla e_K\|_{L^2}^2
+(C_1+C_3)\|e_K\|_{L^2}^2+(C_1+C_3)K^{-2m}\|u\|_{H^m}^2 .
\]
Choosing $\epsilon=D/\bigl(C_{b,\mathrm{lip}}(C_u'+1)\bigr)$ makes the first
bracket $\le0$, so with $C_4:=2(C_1+C_3)$ and $\|e_K\|_{L^2}^2=2\eta$,
\[
\eta'(t)\le C_4\,\eta(t)+C_4\,K^{-2m}\|u\|_{H^m}^2 .
\]
Gr\"onwall's inequality gives
$\eta(t)\le e^{C_4 t}\eta(0)+\int_0^t e^{C_4(t-\sigma)}C_4 K^{-2m}\|u(\sigma)\|_{H^m}^2\intd\sigma$,
hence, with $C_T:=2e^{C_4 T}\max\{1,C_4\}$,
\[
\|e_K(t)\|_{L^2}^2=2\eta(t)\le C_T\Bigl(\|e_K(0)\|_{L^2}^2+K^{-2m}\|u\|_{L^2(0,t;H^m)}^2\Bigr),
\]
which is \eqref{eq:prop-uuK-bound}.
\end{proof}

\begin{proof}[Proof of Proposition~\ref{prop:u_K_mu}]
Recall the birth/death counts of rate $c^\pm=r^\pm(\mathcal P_K\mu^N)$ entering
\eqref{eq:weak-form},
\[
\mathcal R^{\pm,N}_i(t)=Q^{\pm,N}_i\!\Bigl(\int_0^t\mathbf 1_{\mathcal I(N,\tau)}(i)\,
c^\pm(\bX^N_{i,K})\intd\tau\Bigr),\quad
\widetilde{\mathcal R}^{\pm,N}_i(t)=\mathcal R^{\pm,N}_i(t)
-\int_0^t\mathbf 1_{\mathcal I(N,\tau)}(i)\,c^\pm(\bX^N_{i,K}(\tau))\intd\tau,
\]
with $Q^{\pm,N}_i$ independent unit-rate Poisson processes; the net compensator
$c^+-c^-=r$ reproduces the reaction drift in \eqref{eq:weak-form}, and the
predictable quadratic variation of each compensated count is
$\langle\widetilde{\mathcal R}^{\pm,N}_i\rangle_t=\int_0^t\mathbf 1_{\mathcal I(N,\tau)}(i)\,c^\pm(\bX^N_{i,K})\intd\tau$.

Write $\hat e_K=u_K-\mu^N$, $E_K(t)=\tfrac12\|\hat e_K\|_{H^{-s}}^2
=\tfrac12\langle\hat e_K,(I-\Delta)^{-s}\hat e_K\rangle$, and
$\phi_K=(I-\Delta)^{-s}\hat e_K$, so that
\begin{equation}\label{eq:app-phi-id}
\|\phi_K\|_{H^s}^2=\langle(I-\Delta)^{-s}\hat e_K,\hat e_K\rangle=\|\hat e_K\|_{H^{-s}}^2=2E_K(t).
\end{equation}
The process $u_K$ is deterministic and $\mu^N$ is the semimartingale
\eqref{eq:weak-form}. Applying It\^o's formula to the quadratic functional $E_K$,
the time derivative of $\mathbb EE_K$ is the sum of the drift terms and the
predictable quadratic variations of the Brownian and compensated-Poisson
martingales:
\[
\begin{aligned}
\frac{d}{dt}\mathbb EE_K(t)
&=\mathbb E\bigl\langle\nabla\phi_K,\,a(\mathcal P_K\mu^N)\mu^N-a(\mathcal P_K u_K)u_K\bigr\rangle
 +\mathbb E\bigl\langle\Delta\phi_K,\,D(u_K-\mu^N)\bigr\rangle\\
&\quad+\mathbb E\bigl\langle\phi_K,\,r(\mathcal P_K u_K)u_K-r(\mathcal P_K\mu^N)\mu^N\bigr\rangle
 +\mathbb E\Bigl[\frac{D}{N^2}\!\sum_{i\in\mathcal I(N,t)}\!\bigl\|\nabla(I-\Delta)^{-s/2}\delta_{\bX^N_{i,K}}\bigr\|_{L^2}^2\Bigr]\\
&\quad+\mathbb E\Bigl[\frac{1}{2N^2}\!\sum_{i\in\mathcal I(N,t)}\!\bigl\|\delta_{\bX^N_{i,K}}\bigr\|_{H^{-s}}^2\,\bigl|r(\mathcal P_K\mu^N(\bX^N_{i,K}))\bigr|\Bigr]
=\sum_{j=1}^5\mathbb E[I_j].
\end{aligned}
\]
The two quadratic-variation terms $I_4,I_5$ use that a birth or death of particle
$i$ shifts $\hat e_K$ by $\mp\tfrac1N\delta_{\bX^N_{i,K}}$, and that the total
jump rate is $c^++c^-=|r|$.

We use two fixed-$K$ facts.
\emph{(M) Multiplier bound:} a $W^{\lceil s\rceil,\infty}$ function is a pointwise
multiplier on $H^\sigma$ for every $|\sigma|\le s$, with
$\|fg\|_{H^\sigma}\le C_{s,d}\|f\|_{W^{\lceil s\rceil,\infty}}\|g\|_{H^\sigma}$ (by
duality for $\sigma<0$); hence by \eqref{eq:prop2-coeff},
$\|a(\mathcal P_K\mu^N)g\|_{H^\sigma}+\|r(\mathcal P_K\mu^N)g\|_{H^\sigma}+\|u_K g\|_{H^\sigma}\le C_{T,K}\|g\|_{H^\sigma}$
for $|\sigma|\le s$. (An $L^\infty$ bound alone does not suffice for $s>1$, as
$L^\infty$ is not an $H^{s-1}$ multiplier.)
\emph{(I) Inverse estimate:} for $g$ band-limited to $|\mathbf k|_\infty\le K$ and
$\alpha\le\beta$, $\|g\|_{H^\beta}\le C_{T,K}\|g\|_{H^\alpha}$; in particular
$\|\mathcal P_K\hat e_K\|_{L^2}\le C_{T,K}\|\hat e_K\|_{H^{-s}}$ (the projection is an
$H^{-s}$-contraction).

\emph{Term $I_1$ (advection).} Split
$I_1=\langle\nabla\phi_K,a(\mathcal P_K\mu^N)(\mu^N-u_K)\rangle
+\langle\nabla\phi_K,(a(\mathcal P_K\mu^N)-a(\mathcal P_K u_K))u_K\rangle=I_{11}+I_{12}$.
For $I_{11}$, by the duality $\langle H^{s-1},H^{-s+1}\rangle$, (M) with $\sigma=s-1$,
$\|\nabla\phi_K\|_{H^{s-1}}\le\|\phi_K\|_{H^s}$, \eqref{eq:app-phi-id}, and Young,
\[
I_{11}\le\|a(\mathcal P_K\mu^N)\nabla\phi_K\|_{H^{s-1}}\|u_K-\mu^N\|_{H^{-s+1}}
\le\frac{C_{T,K}}{2\epsilon}\|\phi_K\|_{H^s}^2+\frac{C_{T,K}\epsilon}{2}\|\hat e_K\|_{H^{-s+1}}^2
=\frac{C_{T,K}}{\epsilon}E_K+\frac{C_{T,K}\epsilon}{2}\|\hat e_K\|_{H^{-s+1}}^2 .
\]
For $I_{12}$, by (M) for $u_K$, the Lipschitz bound on $a$, and (I),
\[
I_{12}\le\|u_K\nabla\phi_K\|_{H^{s-1}}\bigl\|a(\mathcal P_K\mu^N)-a(\mathcal P_K u_K)\bigr\|_{H^{-s+1}}
\le C_{T,K}\|\phi_K\|_{H^s}\,C_{b,\mathrm{lip}}\|\mathcal P_K\hat e_K\|_{L^2}
\le\frac{C_{T,K}}{2\epsilon}E_K+\frac{C_{T,K}\epsilon}{2}\|\hat e_K\|_{H^{-s+1}}^2 .
\]
Hence $I_1\le\frac{C_{T,K}}{\epsilon}E_K+\frac{C_{T,K}\epsilon}{2}\|\hat e_K\|_{H^{-s+1}}^2$.

\emph{Term $I_2$ (diffusion).} Since $\Delta\phi_K=\phi_K-(I-\Delta)\phi_K$,
\[
I_2=D\langle\Delta\phi_K,\hat e_K\rangle
=D\langle\phi_K,\hat e_K\rangle-D\langle(I-\Delta)^{-s+1}\hat e_K,\hat e_K\rangle
=2D\,E_K-D\|\hat e_K\|_{H^{-s+1}}^2 .
\]

\emph{Term $I_3$ (reaction).} As for $I_1$, with (M) for $u_K$ and
$r(\mathcal P_K\mu^N)$, the Lipschitz bound on $r$, (I), and \eqref{eq:app-phi-id},
\[
\begin{aligned}
I_3&=\langle\phi_K,(r(\mathcal P_K u_K)-r(\mathcal P_K\mu^N))u_K\rangle
+\langle\phi_K,r(\mathcal P_K\mu^N)(u_K-\mu^N)\rangle\\
&\le C_{T,K}\|\phi_K\|_{H^s}\,C_{b,\mathrm{lip}}\|\mathcal P_K\hat e_K\|_{L^2}
+C_{T,K}\|\phi_K\|_{H^s}\|\hat e_K\|_{H^{-s}}
\le C_{T,K}\|\phi_K\|_{H^s}\|\hat e_K\|_{H^{-s}}\le C_{T,K}E_K .
\end{aligned}
\]

\emph{Term $I_4$ (diffusion quadratic variation).} As
$(I-\Delta)^{-s/2}\delta_{\bx}$ has Fourier coefficients
$(1+|\mathbf k|^2)^{-s/2}e^{-i\mathbf k\cdot\bx}$,
\[
\bigl\|\nabla(I-\Delta)^{-s/2}\delta_{\bx}\bigr\|_{L^2}^2
=\sum_{\mathbf k}|\mathbf k|^2(1+|\mathbf k|^2)^{-s}=:C_s'<\infty,
\]
finite precisely because $s>d/2+1$ (the summand decays like $|\mathbf k|^{2-2s}$) and
independent of $\bx$. Hence by $|\mathcal I(N,t)|\le C_N N$,
$I_4=\frac{D}{N^2}\sum_{i\in\mathcal I(N,t)}C_s'\le\frac{D C_N C_s'}{N}=:\frac{C_4}{N}$.

\emph{Term $I_5$ (branching quadratic variation).} Likewise additive: the It\^o
correction contracts each jump $\mp\tfrac1N\delta_{\bX^N_{i,K}}$ with itself, and
\[
\|\delta_{\bx}\|_{H^{-s}}^2=\sum_{\mathbf k}(1+|\mathbf k|^2)^{-s}=:C_s<\infty\qquad(s>d/2),
\]
is independent of $\bx$, so with $\|r\|_{L^\infty}\le C_{b,\mathrm{lip}}$,
$I_5=\frac{1}{2N^2}\sum_{i\in\mathcal I(N,t)}\|\delta_{\bX^N_{i,K}}\|_{H^{-s}}^2\,|r(\mathcal P_K\mu^N(\bX^N_{i,K}))|
\le\frac{C_N C_{b,\mathrm{lip}}C_s}{2N}=:\frac{C_5}{N}$.

\emph{Conclusion.} The terms $I_4,I_5$ are the only ones not proportional to $E_K$
and supply the additive forcing $(C_4+C_5)/N$. Collecting,
\[
\frac{d}{dt}\mathbb EE_K(t)\le\Bigl(\frac{C_{T,K}}{\epsilon}+2D+C_{T,K}\Bigr)\mathbb EE_K(t)
+\frac{C_4+C_5}{N}+\Bigl(\frac{C_{T,K}\epsilon}{2}-D\Bigr)\mathbb E\|\hat e_K\|_{H^{-s+1}}^2 .
\]
Since $D>0$, choosing $\epsilon=D/C_{T,K}$ makes the last term nonpositive, leaving,
with a single $N$-independent constant $C=C_{T,K}$,
\[
\frac{d}{dt}\mathbb EE_K(t)\le C\,\mathbb EE_K(t)+\frac{C}{N},
\]
which is \eqref{eq:prop-uKmu-diff}. Gr\"onwall with constant forcing gives
\eqref{eq:prop-uKmu-bound}; if $\mathbb EE_K(0)\le C_s/N$ then
$\mathbb E\|u_K(t)-\mu^N(t)\|_{H^{-s}}^2\le C_{T,K}/N$ for all $t\le T$.
\end{proof}

\begin{proof}[Proof of Theorem~\ref{thm:abstract_error}]
By the triangle inequality and $(p+q)^2\le 2p^2+2q^2$,
\[
\mathbb E\|u(t_n)-\mu^N(t_n)\|_{H^{-s}}^2
\le 2\|u(t_n)-u_K(t_n)\|_{H^{-s}}^2+2\,\mathbb E\|u_K(t_n)-\mu^N(t_n)\|_{H^{-s}}^2 .
\]
For the first term, $s\ge0$ gives the embedding $\|\cdot\|_{H^{-s}}\le\|\cdot\|_{L^2}$,
and Proposition~\ref{prop:u_u_K} with $e_K(0)=u(0)-u_K(0)=0$ yields
$\|u(t_n)-u_K(t_n)\|_{H^{-s}}^2\le\|u(t_n)-u_K(t_n)\|_{L^2}^2\le C_T K^{-2m}\|u\|_{L^2(0,t_n;H^m)}^2$.
For the second term, the initialization gives
$\mathbb EE_K(0)=\tfrac12\mathbb E\|u(0)-\mu^N(0)\|_{H^{-s}}^2\le C_s/(2N)$, so
Proposition~\ref{prop:u_K_mu} (in the form \eqref{eq:prop-uKmu-bound}) yields
$\mathbb E\|u_K(t_n)-\mu^N(t_n)\|_{H^{-s}}^2=2\,\mathbb EE_K(t_n)\le C_{T,K}/N$. Adding
the two bounds gives \eqref{eq:u-muN-cont}.
\end{proof}

\begin{remark}
The continuous-time analysis above is identical for the Poisson and
minimum-variance kernels (Remark~\ref{rem:rate}). Turning the present estimate into
a sharp, $N$-uniform \emph{fully discrete} bound -- coupling the continuous- and
discrete-time flows on a single probability space by the Harris graphical
construction~\cite{harris1972nearest,bansaye2015stochastic} in the
measure-valued/SDE form of~\cite{fournier2004microscopic,hauray2009wasserstein,chaintron2021propagation},
and extending it to a two-field treatment of the Keller--Segel system up to a
stopping time before concentration -- is carried out in the companion paper
described in Section~\ref{sec:theory}.
\end{remark}


\section{Harris-coupled time-refinement for the splitting rate}
\label{app:harris}
On the manufactured-solution benchmark the first-order time-splitting error is much smaller than the
Monte-Carlo error at every affordable $N$, so it cannot be read off the error against the exact solution
(Figure~\ref{fig:mms}, right, blue). This appendix describes how the splitting rate in Figure~\ref{fig:mms}
(orange) and the mode-resolved agreement in Table~\ref{tab:mms_modes} are measured directly on the branching
scheme.

\paragraph{Spectral splitting-bias field.}
By Proposition~\ref{prop:stochastic_branch} the weighted update and the two branching kernels share the one-step
conditional mean $e^{c(\bx)\tau}$, so the \emph{expected} measure of the scheme evolves, independently of the
kernel, under the deterministic split operator
\begin{equation}\label{equ:Qtau}
Q_\tau = T_\tau\circ M_\tau,\qquad (M_\tau u)(\bx)=e^{r(\bx)\tau}\,u(\bx),
\end{equation}
where $T_\tau$ is the transport (drift--diffusion) sub-step.  In the released particle code the prescribed rate
$r$ is evaluated at the pre-transport position of each particle and the multiplier is applied to its transported
copy, so the mean operator is the first-order splitting $Q_\tau=T_\tau\circ M_\tau$; the spectral solve below
uses the same ordering, so $\varphi$ is exactly the mean of the branching paired difference. Because $\mb b$ and $D$ are constant, the
Euler--Maruyama transport step samples the exact transition law, so $T_\tau$ is the exact semigroup, diagonal in
Fourier, $\widehat{(T_\tau u)}_{\mb k}=e^{-\mathrm i\,\mb b\cdot\mb k\,\tau-D|\mb k|^2\tau}\,\widehat u_{\mb k}$.
Integrating \eqref{equ:Qtau} from $u_0=u_{\mathrm{ex}}(0,\cdot)$ to $T$ by a Fourier pseudo-spectral method (no
particles, hence no Monte-Carlo noise) and forming the band-limited successive difference gives the exact mean
splitting-bias field
\begin{equation}\label{equ:phi}
\varphi(\tau)=\mathcal P_K\big[\,Q_\tau^{\,T/\tau}u_0-Q_{\tau/2}^{\,2T/\tau}u_0\,\big]
=\mathbb E\big[\,u^{N,\tau}-u^{N,\tau/2}\,\big],
\end{equation}
with $\mathcal P_K$ the Fourier projection of Section~\ref{sec:interpolation}. For a first-order method
$\varphi(\tau)=\tfrac12 C\tau+O(\tau^2)$; on this benchmark $\|\varphi(\tau)\|$ halves with each halving of
$\tau$ to machine precision (fitted slope $1.000$), and $C$ matches the analytic splitting commutator to about
two percent.

\paragraph{Lineage-coupled (Harris) paired runs.}
To verify that the \emph{stochastic} branching scheme realizes this mean, we run paired simulations at steps
$\tau$ and $\tau/2$ coupled on a single probability space by a Harris (Ulam--Harris) labelling of the branching
tree~\cite{harris1963branching,harris1972nearest}: every individual carries a persistent random key, and all of
its randomness---transport increment, birth indicator, and the keys of its children---is derived from that key by
namespaced hashing. The two members read the same streams. The coarse transport increment over a window is the
sum of the two fine increments it contains (the same Brownian path on nested grids, so a non-branching particle
has identical positions in both members at the shared times). The coarse birth indicator is calibrated,
$\{U_{2s}<e^{r\tau/2}-1\}\cup\{U_{2s+1}<\theta'\}$ with
$\theta'=(e^{r\tau}-e^{r\tau/2})/(2-e^{r\tau/2})$, so that its marginal birth probability is exactly
$e^{r\tau}-1$, the minimum-variance kernel at step $\tau$ in the pure-birth regime $1<e^{r\tau}<2$; matched
births assign the child the same key in both members. The residual coupling noise $d_s-\varphi(\tau)$ (with
$d_s$ the paired difference defined below) therefore arises only from imperfect coupling---chiefly children born
at a half-step in the fine member, which drift an extra $O(\sqrt\tau)$---and is mean-zero.

\paragraph{Estimators.}
Let $d_s=u^{N,\tau}_s-u^{N,\tau/2}_s$ be the paired difference field for seed $s$, so $\mathbb E[d_s]=\varphi(\tau)$.
Its per-seed magnitude is dominated by the mean-zero de-correlation noise, so a norm such as $\|d_s\|$ (or a
particle-level distance) measures that noise rather than the bias. We instead use the matched-filter projection
onto the known direction $\varphi$,
\begin{equation}\label{equ:rho}
\rho=\frac{\langle d_s,\varphi\rangle}{\langle\varphi,\varphi\rangle},\qquad \mathbb E[\rho]=1,
\end{equation}
averaged over $64$ seeds; the magnitude plotted in Figure~\ref{fig:mms} is $P(\tau)=\bar\rho\,\|\varphi(\tau)\|$,
with standard error from the seed spread. The projection reduces the noise from the full reconstruction space to
a single direction. The mode-resolved version (Table~\ref{tab:mms_modes}) replaces $\varphi$ by each dominant
Fourier basis function $b_k$ and tracks the branching coefficient $a_k(\tau)=\overline{\langle
d_s,b_k\rangle}/\langle b_k,b_k\rangle$ across $\tau$, fitting its log--log slope; the branching coefficients
converge at first order mode by mode, matching the exact spectral rate of $\varphi_k(\tau)$ (slope one). For the
aggregate projection we obtain $\bar\rho=1.07\pm0.07$, $1.14\pm0.09$, and $1.12\pm0.14$ at
$\tau=10^{-1},5\times10^{-2},2.5\times10^{-2}$, all consistent with unity, and $P(\tau)$ scales as $\tau^{0.97}$.
The stochastic branching scheme therefore reproduces the deterministic first-order splitting bias predicted by
the spectral mean-field operator.

\section{Allen--Cahn sanity check}
\label{app:allencahn}
As an additional check on a different nonlinearity, we apply the scalar scheme of Section~\ref{sec:method} to the
Allen--Cahn (AC) equation, a standard phase-separation model with a non-monotone reaction. The reaction is
sign-changing but bounded, so it does not produce the localized growth that motivates branching; we include it
only to confirm that the method reproduces the expected interfacial dynamics. The governing equation is
\begin{equation} \label{eq:ac_pde}
    u_t = D\Delta u + u - u^3, \quad \bx \in \Omega,\ t > 0,
\end{equation}
with diffusion coefficient $D=0.01$ on $\Omega = [0, 2\pi]^2$ with periodic boundary conditions, and initial
condition
\begin{equation} \label{eq:ac_ic}
    u_0(\bx) = \sin^2(x_1) \cos^2(x_2), \quad \bx = (x_1, x_2) \in \Omega.
\end{equation}
We use Fourier projection $K=10$ and $N=400{,}000$ particles up to $T=0.8$. Figure~\ref{fig:ac_solution}
compares the particle solution with a finite-difference reference; the two agree on the phase-separation
dynamics, confirming that the scheme handles a sign-changing reaction correctly.
\begin{figure}[t]
    \centering
    \includegraphics[width=0.8\linewidth]{figure/case_allen_cahn.pdf}
    \caption{Allen--Cahn equation \eqref{eq:ac_pde} at $t=0, 0.2, 0.4, 0.6, 0.8$. Top: finite-difference reference; bottom: the present particle method.}
    \label{fig:ac_solution}
\end{figure}

\section{Coupled-system checks: non-conservative logistic Keller--Segel}
\label{app:coupled}
As an additional coupled-system check we consider a non-conservative Keller--Segel variant with logistic cell
growth,
\begin{equation}
\label{equ:nonlinear}
\begin{aligned}
&u_t - \nabla \cdot \left(\nabla u - \frac{4u}{1 + u^2} \nabla v \right) = u(1 - u), & \quad \bx \in \Omega, t > 0, \\
&v_t - \Delta v = u - v, & \quad \bx \in \Omega, t > 0.
\end{aligned}
\end{equation}
The corresponding well-posedness theory is discussed in~\cite{osaki2002exponential, wang2007classical, wrzosek2004global}.  The initial data are
\[
    \begin{aligned}
        u_0(\bx) =&  \exp\!\left( -(x_1 - 0.5\pi)^2 - (x_2 - 0.5\pi)^2 \right)
        + \exp\!\left( -(x_1 - \pi)^2 - (x_2 - \pi)^2 \right) \\
        &+ \exp\!\left( -(x_1 - 1.5\pi)^2 - (x_2 - 1.5\pi)^2 \right), \\
        v_0(\bx) =&  \exp\!\left( -(x_1 - 0.4\pi)^2 - (x_2 - 0.4\pi)^2 \right)
        + \exp\!\left( -(x_1 - 0.8\pi)^2 - (x_2 - 0.8\pi)^2 \right) \\
        &+ \exp\!\left( -(x_1 - 1.2\pi)^2 - (x_2 - 1.2\pi)^2 \right)
        + \exp\!\left( -(x_1 - 1.6\pi)^2 - (x_2 - 1.6\pi)^2 \right).
    \end{aligned}
\]
Here the cell equation contains genuine birth and death through $u(1-u)$.  The reaction step represents this term by the local rate $f_u/u=1-u$ wherever the reconstructed density is positive.  The run uses periodic reconstruction with $K=10$ and $\tau=10^{-3}$.  Figures~\ref{fig:case3_u} and~\ref{fig:case3_v} compare the particle and finite-difference solutions for $u$ and $v$ at $t=0,0.1,0.2,0.3,0.4$.

\begin{figure}
\centering
\includegraphics[width=0.8\linewidth]{figure/case3_u.pdf}
\caption{Logistic Keller--Segel test: particle and finite-difference approximations of $u$ at $t=0,0.1,0.2,0.3,0.4$; the particle method reproduces the finite-difference cell field under genuine birth/death, validating the non-conservative coupled implementation.}
\label{fig:case3_u}
\end{figure}

\begin{figure}
\centering
\includegraphics[width=0.8\linewidth]{figure/case3_v.pdf}
\caption{Logistic Keller--Segel test: particle and finite-difference approximations of $v$ at $t=0,0.1,0.2,0.3,0.4$; the chemical field likewise agrees with the finite-difference reference.}
\label{fig:case3_v}
\end{figure}

\section{Dense high-dimensional reconstruction check}
\label{app:highdim}
The kinetic Keller--Segel test of Section~\ref{sec:kinetic_ks} is the primary high-dimensional representation test
and avoids reconstructing a dense density.  This appendix is supporting material: it validates the dense
reconstruction's Monte-Carlo scaling alone, against a known reference, using a manufactured solution on
$\mathbb T^d$ ($d=4,6$) for $\partial_t u = D\,\Delta u + r(\bz)\,u$ with the separable, strictly positive exact
solution
\begin{equation}\label{equ:highdim_mms}
u_{\mathrm{ex}}(t,\bz)=e^{\gamma t}\prod_{j=1}^d\big(1+a\cos z_j\big),
\qquad
r(\bz)=\gamma+D\,a\sum_{j=1}^d\frac{\cos z_j}{1+a\cos z_j}.
\end{equation}
We run the minimum-variance branching scheme ($D=0.05$, $a=0.5$, $\gamma=0.2$, $T=0.3$, four seeds) and evaluate
the error by Parseval's identity directly from the empirical Fourier coefficients, so no dense tensor-product grid
is ever assembled.

At the matched cutoff ($k_{\max}=1$, the $3$ modes $\{0,\pm1\}$ per coordinate), the relative $L^2$ error of the
full reconstructed density decays at $N^{-1/2}$: from $0.17$ to $0.029$ in $d=4$ and from $0.44$ to $0.11$ in
$d=6$ as $N$ increases to $6.4\times10^4$ ($3.2\times10^4$ in $d=6$), with fitted slopes $-0.50$ and $-0.49$
(Figure~\ref{fig:highdim_mms}).  The cost is not dimension-free: over-resolving to $k_{\max}=7$ ($15$ modes per
coordinate) raises the mode count from $3^d$ to $15^d$ and inflates the error constant by
$\approx\sqrt{(15/3)^d}=5^{d/2}$ -- about $25\times$ in $d=4$ and $120\times$ in $d=6$, to $0.70$ and above $10$,
at the same $N^{-1/2}$ slope.  This is the variance side of the reconstruction term in
Theorem~\ref{thm:abstract_error}, and the quantitative reason the high-dimensional evidence relies on the
field-coupled kinetic test and grid-free diagnostics rather than a dense product reconstruction.

\begin{figure}[t]
\centering
\includegraphics[width=\linewidth]{figure/highdim_mms_convergence.pdf}
\caption{High-dimensional manufactured-solution density reconstruction.  Relative $L^2$ error of the
reconstructed density versus particle number $N$ for $d=4$ (left) and $d=6$ (right), at the matched
frequency cutoff $k_{\max}=1$ (green, $3$ modes per coordinate) and the over-resolved $k_{\max}=7$ (red, $15$ modes
per coordinate).  Both converge at the Monte-Carlo rate $N^{-1/2}$ (fitted slopes shown); the over-resolved error
is larger by a factor that grows as $\sqrt{(2k_{\max}+1)^d}$, illustrating that dense Fourier reconstruction is
convergent but not dimension-free.  Bars are one standard deviation over four seeds.}
\label{fig:highdim_mms}
\end{figure}

The reconstruction also recovers the right \emph{shape}, not only the global rate.  For a $d=6$ run at
$N=6.4\times10^4$ (cutoff $k_{\max}=1$), all six one-dimensional coordinate marginals coincide with the exact
normalized marginal $(1+a\cos z)/2\pi$, and the two-dimensional marginals of $(z_1,z_2)$ are visually
indistinguishable from the exact ones (Figure~\ref{fig:highdim_mms_marg}): the mean one-dimensional $L^1$ marginal
error is $5.0\times10^{-3}$ and the two-dimensional relative $L^2$ marginal error is $1.1\times10^{-2}$
($4.0\times10^{-3}$ and $8.0\times10^{-3}$ in $d=4$).

\begin{figure}[t]
\centering
\includegraphics[width=\linewidth]{figure/highdim_mms_marginals_d6.pdf}
\caption{Marginals of the reconstructed density for the $d=6$ manufactured-solution run ($N=6.4\times10^4$, cutoff
$k_{\max}=1$).  Left: the reconstructed one-dimensional marginal of each of the six coordinates (grey) against the
exact marginal $(1+a\cos z)/2\pi$ (dashed).  Middle and right: the reconstructed and exact two-dimensional
marginals of $(z_1,z_2)$, on a common colour scale.}
\label{fig:highdim_mms_marg}
\end{figure}

\printcredits

\bibliographystyle{elsarticle-num}
\bibliography{main.bib}
\end{document}